% Options for packages loaded elsewhere
\PassOptionsToPackage{unicode}{hyperref}
\PassOptionsToPackage{hyphens}{url}
\PassOptionsToPackage{dvipsnames,svgnames,x11names}{xcolor}
\documentclass[
  a4paper,
  11pt]{article}
\usepackage{xcolor}
\usepackage{amsmath,amssymb}
\setcounter{secnumdepth}{5}
\usepackage{iftex}
\ifPDFTeX
  \usepackage[T1]{fontenc}
  \usepackage[utf8]{inputenc}
  \usepackage{textcomp} % provide euro and other symbols
\else % if luatex or xetex
  \usepackage{unicode-math} % this also loads fontspec
  \defaultfontfeatures{Scale=MatchLowercase}
  \defaultfontfeatures[\rmfamily]{Ligatures=TeX,Scale=1}
\fi
\usepackage{lmodern}
\ifPDFTeX\else
  % xetex/luatex font selection
\fi
% Use upquote if available, for straight quotes in verbatim environments
\IfFileExists{upquote.sty}{\usepackage{upquote}}{}
\IfFileExists{microtype.sty}{% use microtype if available
  \usepackage[]{microtype}
  \UseMicrotypeSet[protrusion]{basicmath} % disable protrusion for tt fonts
}{}
\makeatletter
\@ifundefined{KOMAClassName}{% if non-KOMA class
  \IfFileExists{parskip.sty}{%
    \usepackage{parskip}
  }{% else
    \setlength{\parindent}{0pt}
    \setlength{\parskip}{6pt plus 2pt minus 1pt}}
}{% if KOMA class
  \KOMAoptions{parskip=half}}
\makeatother
\setlength{\emergencystretch}{3em} % prevent overfull lines
\providecommand{\tightlist}{%
  \setlength{\itemsep}{0pt}\setlength{\parskip}{0pt}}
% Preamble for the AETHER paper. Loaded by pandoc -H; runs under XeTeX (tectonic).
%
% Type: Charter (Matthew Carter, 1987) for the text, Avenir Next for headings, captions,
% table heads and running heads — the same sans the paper's infographics are set in, so
% figures and page share one type system — Menlo for run IDs and file paths only, and
% STIX Two Math as the glyph source for the few symbols Charter lacks (primes, arrows).
% Palette: the infographics' own ink / rule / accent-blue tokens (scripts/render_infographics.py).

\usepackage[a4paper, left=28mm, right=28mm, top=26mm, bottom=28mm, headsep=9mm, footskip=13mm]{geometry}

\usepackage{fontspec}
\setmainfont{Charter}[
  BoldFont = {Charter Bold}, ItalicFont = {Charter Italic}, BoldItalicFont = {Charter Bold Italic}]
\setsansfont{Avenir Next}[
  UprightFont = {Avenir Next Regular}, BoldFont = {Avenir Next Demi Bold},
  ItalicFont = {Avenir Next Italic}, BoldItalicFont = {Avenir Next Demi Bold Italic}]
\setmonofont{Menlo}[Scale=0.84]
\newfontfamily\symfont{STIX Two Math}
\newfontfamily\sfmedium{Avenir Next Medium}[Ligatures=TeX]

\usepackage{xcolor}
\definecolor{ink}{HTML}{141413}
\definecolor{inktwo}{HTML}{52514e}
\definecolor{muted}{HTML}{898781}
\definecolor{ruleline}{HTML}{c3c2b7}
\definecolor{panel}{HTML}{f4f4f1}
\definecolor{accent}{HTML}{184f95}
\definecolor{good}{HTML}{0ca30c}
\definecolor{warn}{HTML}{e09a00}
\definecolor{crit}{HTML}{d03b3b}
\color{ink}

% reference-tier marks (the markdown uses emoji; XeTeX cannot set colour emoji)
\usepackage{pifont}
\newcommand{\tierok}{{\color{good}\ding{51}}}
\newcommand{\tiermid}{{\color{warn}\ding{108}}}
\newcommand{\tierno}{{\color{crit}\ding{55}}}

% body rhythm
\linespread{1.10}
\setlength{\parindent}{0pt}
\setlength{\parskip}{5pt plus 1pt minus 1pt}
\emergencystretch 2.5em
\setcounter{secnumdepth}{2}
\setcounter{tocdepth}{1}
\clubpenalty=8000 \widowpenalty=8000

% headings
\usepackage[nobottomtitles*]{titlesec}
\renewcommand{\bottomtitlespace}{0.12\textheight}
\titleformat{\section}
  {\sffamily\bfseries\fontsize{15.5}{19}\selectfont\color{ink}}
  {{\color{accent}\thesection}}{0.75em}{}
\titlespacing*{\section}{0pt}{24pt plus 5pt minus 3pt}{9pt}
\titleformat{\subsection}
  {\sffamily\bfseries\fontsize{12}{15}\selectfont\color{ink}}
  {{\color{accent}\thesubsection}}{0.7em}{}
\titlespacing*{\subsection}{0pt}{15pt plus 3pt minus 2pt}{5pt}
\titleformat{\subsubsection}
  {\sfmedium\fontsize{10.5}{13}\selectfont\color{inktwo}}
  {}{0em}{}
\titlespacing*{\subsubsection}{0pt}{11pt plus 2pt minus 1pt}{3pt}
\usepackage{placeins}   % \FloatBarrier before the references and the appendix only

% running heads
\usepackage{fancyhdr}
\pagestyle{fancy}
\fancyhf{}
\fancyhead[L]{\sffamily\fontsize{8}{10}\selectfont\color{inktwo}AETHER\;·\;When Cooler Is Not Safer}
\fancyhead[R]{\sffamily\fontsize{8}{10}\selectfont\color{inktwo}Draft, 21 September 2026 · AI-drafted, not yet revised by the author}
\fancyfoot[C]{\sffamily\fontsize{9}{11}\selectfont\color{inktwo}\thepage}
\renewcommand{\headrulewidth}{0pt}
\fancypagestyle{plain}{\fancyhf{}\fancyfoot[C]{\sffamily\fontsize{9}{11}\selectfont\color{inktwo}\thepage}\renewcommand{\headrulewidth}{0pt}}

% captions: sans, small, label in the accent
\usepackage{caption}
\captionsetup{font={sf,footnotesize}, labelfont={bf,color=accent}, labelsep=period,
  justification=justified, singlelinecheck=false, skip=7pt}
\captionsetup[table]{position=above, skip=5pt}

% floats: let a figure page hold a tall figure, keep text with small ones
\renewcommand{\topfraction}{0.92}
\renewcommand{\bottomfraction}{0.6}
\renewcommand{\textfraction}{0.06}
\renewcommand{\floatpagefraction}{0.7}
\setcounter{topnumber}{3}
\setcounter{totalnumber}{4}
\makeatletter   % on a float-only page, stack floats from the top instead of centring them
\setlength{\@fptop}{0pt}
\setlength{\@fpsep}{14pt plus 2pt}
\setlength{\@fpbot}{0pt plus 1fil}
\makeatother

% boxes: the page-1 status banner and block quotes (hypotheses, the pending-results slot)
\usepackage{tcolorbox}
\tcbuselibrary{breakable, skins}
\newtcolorbox{draftbanner}{enhanced, breakable, colback=panel, colframe=panel,
  borderline west={2.5pt}{0pt}{accent}, arc=0pt, boxrule=0pt, boxsep=0pt,
  left=10pt, right=10pt, top=6pt, bottom=6pt, before skip=8pt, after skip=10pt,
  fontupper=\sffamily\fontsize{9}{11.6}\selectfont\color{ink}}
\renewenvironment{quote}
  {\begin{tcolorbox}[enhanced, breakable, colback=white, colframe=white,
     borderline west={1.5pt}{0pt}{ruleline}, arc=0pt, boxrule=0pt, boxsep=0pt,
     left=10pt, right=0pt, top=3pt, bottom=3pt, before skip=6pt, after skip=8pt]}
  {\end{tcolorbox}}

% tables
\usepackage{booktabs, longtable, array, colortbl}
\setlength{\LTpre}{6pt}
\setlength{\LTpost}{10pt}
\setlength{\tabcolsep}{5pt}
\renewcommand{\arraystretch}{1.18}
\arrayrulecolor{ruleline}
\newcommand{\tabhead}[1]{{\sffamily\bfseries\fontsize{8.5}{10.5}\selectfont #1}}
\setlength{\aboverulesep}{0.3ex}
\setlength{\belowrulesep}{0.5ex}

% contents page
\usepackage{tocloft}
\renewcommand{\cfttoctitlefont}{\sffamily\bfseries\fontsize{15.5}{19}\selectfont}
\renewcommand{\cftaftertoctitle}{\vspace{4pt}}
\renewcommand{\cftsecfont}{\sffamily\bfseries\small}
\renewcommand{\cftsecpagefont}{\sffamily\small}
\renewcommand{\cftsubsecfont}{\small}
\renewcommand{\cftsubsecpagefont}{\small}
\renewcommand{\cftsecleader}{\cftdotfill{\cftsecdotsep}}
\renewcommand{\cftsecdotsep}{\cftnodots}
\renewcommand{\cftsubsecdotsep}{4.5}
\setlength{\cftbeforesecskip}{5pt}
\setlength{\cftbeforesubsecskip}{1pt}
\setlength{\cftsecnumwidth}{2.2em}
\setlength{\cftsubsecindent}{2.2em}
\setlength{\cftsubsecnumwidth}{2.6em}

% lists a little tighter
\usepackage{enumitem}
\setlist{itemsep=2pt, parsep=0pt, topsep=4pt, leftmargin=1.6em}
\setlist[enumerate]{leftmargin=2.2em}
\usepackage{bookmark}
\IfFileExists{xurl.sty}{\usepackage{xurl}}{} % add URL line breaks if available
\urlstyle{same}
\hypersetup{
  colorlinks=true,
  linkcolor={accent},
  filecolor={Maroon},
  citecolor={Blue},
  urlcolor={accent},
  pdfcreator={LaTeX via pandoc}}

\author{}
\date{}

\begin{document}

% Title block, page 1. Title and author line as they stand in AETHER_paper.md and README.md;
% technical subtitle verbatim from CLAUDE.md §1. No affiliation: the paper is an independent capstone.
\thispagestyle{plain}
\enlargethispage{2\baselineskip}
\noindent
{\sffamily\bfseries\fontsize{9}{11}\selectfont\color{accent}AETHER}%
{\sffamily\fontsize{9}{11}\selectfont\color{inktwo}\;·\;independent capstone\;·\;2026\par}
\vspace{6pt}
{\sffamily\bfseries\fontsize{22}{26}\selectfont\color{ink}\raggedright
When Cooler Is Not Safer:\\
Peak and In-Depth Thermal Optimisation of Atmospheric Re-entry\par}
\vspace{7pt}
{\rmfamily\itshape\fontsize{12.5}{16}\selectfont\color{inktwo}\raggedright
AI-Guided Multifidelity Framework for Coupled Re-entry Trajectory, Aerodynamics and TPS Optimization\par}
\vspace{11pt}
{\sffamily\fontsize{11.5}{14}\selectfont\color{ink}Adithya Kesan Jayakanth\par}
\vspace{3pt}
{\sffamily\fontsize{9}{11}\selectfont\color{inktwo}Draft of 21 September 2026\;·\;typeset from repository commit \texttt{de2a834}\par}
\vspace{7pt}
{\color{ruleline}\rule{\linewidth}{0.6pt}\par}
\vspace{2pt}

\begin{draftbanner}
Draft status, 2026-09-21. This text was drafted by an AI assistant from
the repository's result files at commit \texttt{de2a834} and has not yet
been revised by the author. Every number is read from a file named in
the text, a table note or a caption. The author must revise it, check it
against the files, and own it before it is submitted anywhere. See §19.
\end{draftbanner}

\section{Abstract}

A re-entry heat shield is normally designed against peak heat flux, the
worst instant of heating on its outer surface, but it does not fail
there. It fails at the bondline, where the thermal protection meets the
structure, and the bondline responds to the heat that gets past the
surface over the whole entry. This work asks whether the two quantities
are minimised by the same entry, and whether optimising both together
finds designs that optimising the peak alone would miss.

A reduced-order chain (US Standard Atmosphere 1976, point-mass entry,
Sutton--Graves stagnation heating with an effective nose radius taken
from measured stagnation-point velocity gradients, and one-dimensional
transient conduction with a radiating surface) was checked against
published references before use. The conduction solver agrees with an
analytical solution to 0.002\%, and the trajectory integrator reproduces
a published closed-form-versus-numerical comparison to within 0.92
percentage points. Drag comes from a response surface fitted to 69
inviscid perfect-gas CFD cases.

Over a 41-point sweep of entry angle, peak heat flux and peak bondline
temperature are strictly monotone with opposite signs. On the final
model, the knee of the two-objective Pareto front runs 14.9 K cooler at
the bondline than the peak-flux-only optimum, for 21.5
kW/m\textsuperscript{2} more peak flux; under the declared uncertainties
the paired difference is −14.31 K (s.d. 1.45 K) and is negative in all
3000 draws. That result is a statement about entry steepness at a
geometry fixed by placeholder constraints, not about a located capsule
shape. The hypothesis that adaptive use of CFD would reduce the number
of CFD runs needed was not supported. The physical coupon experiment has
not been run.

All results are model predictions within a documented set of
assumptions; no flight or experimental validation is claimed.

\tableofcontents\clearpage

\section{Introduction}

A thermal protection system (TPS) is sized against two different
numbers, and the design literature is explicit about which does which
job. NASA's aerothermodynamics training material states that heat flux,
with pressure and shear, selects the TPS material, while heat load, the
time integral of that flux, determines its thickness {[}2{]}. The same
source states the consequence for trajectory design as a known trade:
``The selection of γ becomes a trade between peak heat rate (TPS
material selection), and total heat load (TPS thickness and mass)''
{[}2{]}. A steeper entry gives a higher peak; a shallower one gives a
higher total.

So the trade is known. What is stated less often is where the failure it
guards against occurs. The outer surface of a heat shield is designed to
get very hot and to radiate most of the incoming energy back out. The
plane that matters structurally is the bondline, and the accepted
in-depth criterion is a maximum temperature at that single plane, held
below an allowable that comes from the structure rather than from the
insulator {[}2{]}.

That allowable differs between flown vehicles. The Space Shuttle
Orbiter's aluminium primary structure carried a design limit of 350 °F,
449.8 K {[}12{]}. The Apollo Command Module's
ablator-to-stainless-honeycomb interface was held to 600 °F, 588.7 K
{[}11{]}. Orion's Avcoat over a composite skin on a titanium skeleton is
quoted against 500 °F, 533.2 K {[}13{]}. In each case the number is set
by what is behind the heat shield.

Peak heat flux is an instantaneous surface quantity. Bondline
temperature is what remains after the whole pulse has been filtered by
conduction through the stack, and it peaks late, after aeroheating has
stopped, because heat already inside the material keeps diffusing
inward. Stagnation heating goes as V\textsuperscript{3}√ρ, so a
shallower entry, which decelerates higher in thinner air, has a lower
peak. It also lasts much longer. Whether that matters depends on one
comparison: the diffusion time of the stack, L\textsuperscript{2}/α,
against the duration of the entry. If the stack is much thicker than the
diffusion length √(αt), the bondline never registers the entry. If it is
much thinner, the bondline tracks the surface. When the two timescales
are comparable, the duration of the pulse changes the answer as well as
its size. A heat shield with no spare insulation sits in that regime,
because spare insulation is mass.

\begin{figure}[htbp]
\centering
\includegraphics[width=\linewidth,height=0.42\textheight,keepaspectratio]{/Users/akj/Projects/aether/reports/final/../figures/infographics/01_burn_vs_bake_mechanism.pdf}
\caption{Burn versus bake, the mechanism: a shallow entry lowers peak heat flux
by 38.3\% and raises the peak bondline temperature by 114.1 K (M1,
Fidelity 0, run \texttt{M1-\allowbreak{}20260902T130547Z}). Explanatory
figure: the heat-flux and in-depth histories are re-derived through
\texttt{evaluate\_\allowbreak{}design()} from the run's configuration
snapshot and the scalars are asserted equal to \texttt{candidates.csv};
legacy cap-radius nose model, as the M1 run used. File
\texttt{reports/\allowbreak{}figures/\allowbreak{}infographics/\allowbreak{}01\_\allowbreak{}burn\_\allowbreak{}vs\_\allowbreak{}bake\_\allowbreak{}mechanism.png}.}
\end{figure}

This project asks what follows. Can minimising peak external heat flux
alone select a trajectory that is worse at the bondline? Does optimising
the peak and the in-depth response jointly find designs a peak-only
procedure would miss? The objective is kept as a multi-objective problem
throughout and is never collapsed into a weighted score.

The approach was to build a reduced-order chain, check each component
against a reference outside the project before using it, demonstrate the
effect where the model is defensible, and then add a CFD-derived drag
model under an explicit gate: no CFD result could enter the optimisation
until mesh independence, force convergence and a published blunt-body
comparison were documented. Optimisation, an optimiser comparison, an
adaptive-fidelity study and an uncertainty study followed on that model.

The claim is narrow on purpose. Within a verified reduced-order model
and a documented set of assumptions, peak heat flux and peak bondline
temperature are optimised by different entries, and the knee of the
joint front has a cooler bondline than the peak-flux-only optimum by an
amount that survives the declared uncertainties. The paper claims
nothing about a flight vehicle, a qualified material or a flight
condition, and no experimental validation exists.

One thing about provenance belongs here and not in a footnote. The
research question, the locked objective and the three hypotheses were
set out in a project specification that I received and accepted. The
framing that minimising peak heat flux may not give the thermally safest
entry came from that specification, not from an observation made during
this work. What was done here is the modelling, the verification, the
attempts to falsify, and the decisions about what the evidence supports.
The dated record of how the question developed is
\texttt{docs/\allowbreak{}research\_\allowbreak{}story.md}.

\section{Related Work}

Every citation in this section was opened and read during the project
and recorded in
\texttt{docs/\allowbreak{}validation/\allowbreak{}sourcing\_\allowbreak{}report.md}
or \texttt{docs/\allowbreak{}theory/\allowbreak{}*}, or is marked as
secondary. Sources that could not be obtained are named as gaps. Tiers
appear in the reference list.

\subsection{Stagnation-point heating correlations}

Sutton and Graves fitted a stagnation-point heating coefficient to
equilibrium boundary-layer solutions across nine base gases and
twenty-two mixtures {[}1{]}. Their equation is written in stagnation
pressure and enthalpy, \.{q}\_w = K √(p\_s/R)(h\_s − h\_w), with K =
0.1113
kg·s\textsuperscript{−1}·m\textsuperscript{−3/2}·atm\textsuperscript{−1/2}
for air, an average correlation error of 3.3\% and a maximum of 9.8\%,
fitted over stagnation enthalpies of 2.3 to 116.2 MJ/kg and pressures of
0.001 to 100 atm. The report states that radiative heating was
neglected.

The form used throughout the entry literature, \.{q}{\symfont ″} = k
√(ρ\_∞/R\_n) V\_∞\textsuperscript{3}, does not appear in that report,
and neither does the constant usually attributed to it. Reaching it
needs three substitutions: a Newtonian stagnation pressure, a total
enthalpy of V\textsuperscript{2}/2, and a cold wall. The chain is given
in NASA training material {[}2{]} and in Carroll and Brandis {[}3{]},
both secondary. Section 6.3 reports what happened when I carried the
algebra through from the primary.

The largest unquantified term in this class of correlation is surface
catalycity. Secondary sources characterise the
equilibrium-boundary-layer basis as equivalent to a fully catalytic
wall, and report catalytic heat flux at roughly twice non-catalytic
values for comparable conditions {[}4{]}. This project does not quantify
it.

The standard critical review of these correlations is attributed to
Tauber. The repository's own records of whether that document was read
contradict each other, so no claim in this paper rests on it (see the
second block of §20).

\subsection{Stagnation-point velocity gradients on blunt bodies}

The 1/√R\_n scaling is a statement about the stagnation-point velocity
gradient, and it is exact only for a sphere. Zoby and Sullivan treat the
general blunt body by defining an effective radius: the radius of the
hemisphere whose stagnation velocity gradient equals the body's,
R\_b/R\_eff = (dU/dS)\_BB/(dU/dS)\_hemi, their equation (5) {[}5{]}.
Their method rests on the stagnation-region pressure distribution being
invariant above about Mach 3.5. Ellison measured the same quantity on
nine models at Mach 8, tabulating R\_b/R\_eff over K = R\_b/R\_n in \{0,
0.417, 0.707\} and corner ratio R\_c/R\_b in \{0, 0.2, 0.4\} {[}6{]}.
Two features of that table matter here. The effective radius saturates:
a flat face with a sharp corner reaches 3.155 R\_b, not infinity. And a
rounder corner raises stagnation heating.

Ellison states the disagreement between the two sources himself: within
10\% at K = 0 and K = 0.707, and about 20\% at K = 0.417 with a sharp
corner {[}6{]}. K = 0.417 is where this project's optimised designs sit.
Substituting an effective radius into a Sutton--Graves-form correlation
has a published precedent in Zoby's own later work {[}7{]}.

Neither source publishes a formula. The interpolation used here is this
project's construction over their data points. A 2026 paper whose title
describes a closed-form reconstruction of the Zoby--Sullivan scaling
exists and is paywalled; it has not been read.

\subsection{Ballistic entry trajectory solutions}

Allen and Eggers give the classical closed-form treatment under a
straight flight path, negligible gravity, an exponential atmosphere and
constant drag coefficient {[}8{]}. Because the closed form neglects
gravity, a correct numerical integrator must disagree with it, more so
as entry flattens. Putnam and Braun quantify that disagreement for three
entry cases, with peak-deceleration differences of +34.9\%, −5.1\% and
−57.3\% for sample-return, steep-strategic and shallow LEO-return
entries {[}9{]}. That published disagreement is the validation target
used in §6.2. The Stardust reconstruction gives a flight value, 32.89 g
measured against 32.86 g predicted {[}10{]}, but comparing against it
needs a drag coefficient for that capsule that could not be sourced. Two
standard textbooks on entry mechanics (Vinh, Busemann and Culp; Regan
and Anandakrishnan) could not be obtained.

\subsection{TPS sizing criteria and limits}

NASA training material describes baseline TPS sizing as computed to a
bondline temperature limit, with thermal margin applied to that
criterion {[}2{]}. The three flown limits of §2 are all read from
primaries {[}11--13{]}. One qualification: CR-159900 states 350 °F as
the aluminium airframe's design limit and rates the RTV-560 adhesive
itself to 500 °F; no primary sentence was found tying 350 °F to the
adhesive line.

For deceleration, no document states a flat scalar. NASA-STD-3001 gives
a duration-dependent sustained-acceleration curve; for a deconditioned
crew it runs from 14.0 g at 0.5 s to 4.0 g beyond 150 s {[}14{]}.
Published peak decelerations of crewed entries are 7.6 to 11.1 g for
Mercury, 4.3 to 7.7 g for Gemini, 3.3 to 6.8 g for Apollo and 3 to 4 g
for nominal Soyuz {[}15{]}.

\subsection{Cumulative thermal-exposure metrics}

The project defined and tested a Thermal Penetration Index (TPI), a
weighted double integral of temperature exceedance over depth and time.
Such metrics have several independent literatures: integrated heat load
in aerothermal design {[}2{]}; the CEM43 thermal dose in hyperthermia,
read through a consensus paper that reproduces it {[}16{]}; and the
Arrhenius damage integral of burn-injury work, read through a review
{[}17{]}. A search for a named metric of the TPI's exact shape in TPS
design found none in the sources it could reach. That is a negative
search result, not a proof of absence and not a novelty claim;
paywalled, industry-internal and export-controlled material was not
reachable.

\subsection{Base drag, and flown capsule shapes}

The CFD here computes forebody drag only, so base drag comes from a
stated assumption. Chapman {[}18{]} and Love {[}19{]} give the
semi-empirical foundations; Miller's wind-tunnel study of blunted 9°
cones at Mach 10.5 to 20 reports base drag under about 2\% of total for
bluntness ratios of 0.55 to 0.8 {[}20{]}; Hoerner gives the qualitative
trend {[}21{]}. Every quantitative dataset found is low-Reynolds-number
tunnel data. Flight backshell pressure from Mars 2020 exists and was
reached at abstract level only {[}22{]}.

Nine flown vehicles were checked in primaries to place the design box
{[}23{]}. Shallow spherical-cap capsules cluster at R\_n/D of about 1.0
to 1.2 (Apollo 1.20, Mercury 1.07); blunted sphere-cones cluster at 0.25
to 0.5. Apollo's commonly quoted 32.5° half-angle could not be confirmed
in a primary, whose baseline configuration gives 33.0°.

No literature on multi-objective optimisation or on machine-learning
surrogates was opened for this project, so none is cited.

\section{Hypotheses}

Three hypotheses were frozen in the specification before any model was
built, and are reproduced verbatim.

\begin{quote}
\textbf{H0.} Minimising peak heat flux alone does not necessarily
produce the thermally safest re-entry solution.

\textbf{H1.} Joint optimisation of peak heat flux and in-depth TPS
response can identify safer feasible designs than peak-heat-only
optimisation.

\textbf{H2.} Surrogate-assisted and AI-guided adaptive-fidelity search
can approach the Pareto frontier with fewer expensive CFD evaluations
than unguided or brute-force search.
\end{quote}

H0 is an existence claim. It asks for one pair of designs A and B with
q{\symfont ″}\_max(B) \textless{} q{\symfont ″}\_max(A) and
T\_bond,max(B) \textgreater{} T\_bond,max(A) inside the model's valid
domain, and it is falsified by searching and finding none. The search is
automated, and a pair counts only if the flux reduction is at least
1.0\% and the bondline penalty at least 1.0 K; both thresholds were
declared in advance. The falsifying outcome is reachable and was reached
once: with a 40 mm insulator the bondline rose 8.8 K over a whole entry,
so for that stack the two metrics are effectively independent (NR-03).
That is what makes the later positive result a statement about a regime.

H1 is a comparison, falsified if a joint optimiser and a peak-flux-only
optimiser over the same feasible set select designs whose bondline
temperatures do not differ beyond the model's resolution. Two
definitions were fixed first. ``Joint'' means the knee of the combined
feasible front, the point furthest from the chord joining the front's
extremes in normalised objective space, which contains no weights.
``Safer'' is reported as an absolute temperature difference, because a
margin ratio swings from 37.5× to 1.6× as the bondline allowable moves
from 445 K to 500 K.

H2 was given a criterion in
\texttt{configs/\allowbreak{}adaptive\_\allowbreak{}fidelity.yaml}
before any run. Cost is the number of CFD calls charged when an arm's
truth hypervolume first reaches 95\% of a reference value. H2 is
\texttt{SUPPORTED} only if the arm that buys no CFD fails to reach the
target in at least three of five seeds, the adaptive arm does reach it,
its median cost is below the best non-adaptive arm's, and a one-sided
exact permutation test gives p \textless{} 0.05. If the no-CFD arm
reaches the target the verdict is \texttt{NOT\_\allowbreak{}TESTABLE},
since ``fewer calls'' means nothing if none were needed.

\begin{center}
\small\centering
\begin{tabular}{@{}>{\raggedright\arraybackslash}p{\dimexpr(\linewidth-8\tabcolsep)*98/1000\relax}>{\raggedright\arraybackslash}p{\dimexpr(\linewidth-8\tabcolsep)*242/1000\relax}>{\raggedright\arraybackslash}p{\dimexpr(\linewidth-8\tabcolsep)*340/1000\relax}>{\raggedright\arraybackslash}p{\dimexpr(\linewidth-8\tabcolsep)*318/1000\relax}@{}}
\toprule
\tabhead{Hypothesis} & \tabhead{Status} & \tabhead{Evidence} & \tabhead{Model} \\
\midrule
H0 & Supported in the tested domain & §7;
\texttt{reports/\allowbreak{}milestones/\allowbreak{}M1\_\allowbreak{}burn\_\allowbreak{}vs\_\allowbreak{}bake.md},
\texttt{M1\_\allowbreak{}signature\_\allowbreak{}counterexample.md};
restated under uncertainty in §13 & Fidelity 0 (constant C\_D), then Fidelity 1 \\
H1 & Supported, narrowly (see §10.4) & §7.4 (Fidelity 0); §10 and §13 (Fidelity 1) & both \\
H2 & Not supported & §12;
\texttt{reports/\allowbreak{}milestones/\allowbreak{}M6\_\allowbreak{}adaptive\_\allowbreak{}fidelity.md} & Fidelity 1 \\
\bottomrule
\end{tabular}
\end{center}

None of the three says anything about a real vehicle or a qualified
material. H0 and H1 are statements about a model. H2 is a statement
about search efficiency in that model. Whether the conduction physics
behaves as modelled in a real solid is the business of the coupon
experiment, which has not been run (§14).

\section{Methods}

Two model versions appear in this paper. Fidelity 0 has a constant drag
coefficient of 1.20, so forebody shape has no aerodynamic consequence
and only diameter acts, through reference area. Fidelity 1 replaces the
constant with the CFD-derived drag surface of §9. Every result is
labelled with the one it came from.

\subsection{The canonical evaluator}

One function, \texttt{evaluate\_\allowbreak{}design(config)}, runs the
whole chain: validate the geometry, obtain the aerodynamic model,
integrate the trajectory, compute the heat-flux history, solve the TPS
response, compute objectives, check constraints, attach provenance,
return a status. No study bypasses it. The uncertainty propagation,
which would naturally loop over draws outside the optimiser, was
implemented by appending the draw index to the design space as one
synthetic variable, so each design-and-draw pair is an ordinary design
vector that is charged and logged like any other (A-UQ;
\texttt{uncertainty/\allowbreak{}space.py}). Three runs of the evaluator
on one configuration produce identical result fingerprints covering
every metric, margin, diagnostic, the provenance block and the velocity
history
(\texttt{reports/\allowbreak{}milestones/\allowbreak{}M3\_\allowbreak{}coupled\_\allowbreak{}model.md}
§10).

\begin{figure}[htbp]
\centering
\includegraphics[width=\linewidth,height=0.42\textheight,keepaspectratio]{/Users/akj/Projects/aether/reports/final/../figures/infographics/02_pipeline_validation.pdf}
\caption{The pipeline and what is validated against what. Gate statuses are the
validation matrix's own (\texttt{VALIDATION\_\allowbreak{}MATRIX.md},
2026-09-21); G4 \texttt{PASS} carries its restart history (§8.2). Box
descriptions paraphrase the matrix's validation-source column; no number
on the figure is typed by hand. Sources:
\texttt{results/\allowbreak{}M2/\allowbreak{}M2-\allowbreak{}20260920T123901Z/\allowbreak{}gate\_\allowbreak{}assessment.json},
\texttt{results/\allowbreak{}M4/\allowbreak{}M4-\allowbreak{}OPT-\allowbreak{}20260920T205252Z/\allowbreak{}summary.json},
\texttt{results/\allowbreak{}M7/\allowbreak{}M7-\allowbreak{}UQ-\allowbreak{}20260920T233048Z/\allowbreak{}summary.json},
\texttt{ARCHITECTURE.md}. File
\texttt{reports/\allowbreak{}figures/\allowbreak{}infographics/\allowbreak{}02\_\allowbreak{}pipeline\_\allowbreak{}validation.png}.}
\end{figure}

\subsection{Atmosphere}

From 0 to 86 km the US Standard Atmosphere 1976 layer profile is
integrated exactly in geopotential altitude with an ideal-gas closure
{[}24{]}. From 86 to 150 km density is log-interpolated from a
transcribed table, a stand-in for the species-diffusion model the
standard defines there, and every value from that region is flagged
\texttt{extrapolated}. Above 150 km the model raises an error (A-ATM-1
to A-ATM-3).

The assumption first attached to the upper region was that its
contribution is negligible. That is true for peak heat flux and false
for the bondline (§17.1).

\subsection{Trajectory}

A point-mass, three-degree-of-freedom planar entry over a non-rotating
spherical Earth, state {[}altitude, velocity, flight-path angle,
range{]}, integrated with an adaptive Runge--Kutta scheme and terminal
events; inverse-square gravity; Earth radius 6371 km; integration ends
at 20 km and the thermal solve continues past that (A-TRAJ-1 to
A-TRAJ-5). Earth rotation would change relative velocity by up to about
0.46 km/s at the equator, which is up to about 19\% on heat flux; it
shifts magnitudes and does not reverse the ordering between steep and
shallow entries. There is no lift, bank or angle of attack, and mass is
constant. An ablator would carry energy away as mass loss, more in the
long shallow case, which would reduce the reported effect.

\subsection{Aeroheating}

Stagnation-point convective heating, \.{q}{\symfont ″} = k
√(ρ\_∞/R\_eff) V\_∞\textsuperscript{3} with k =
1.7415×10\textsuperscript{−4} kg\^{}½ m\textsuperscript{−1}, and the
integrated external heat load as its time integral (A-HEAT-1 to
A-HEAT-4). It is convective only and cold-wall. It is evaluated at the
stagnation point only: the model says nothing about the shoulder or
afterbody, which is a limitation and not a conservatism, and §10 shows
an optimiser finding it.

\subsubsection{The effective nose
radius}\label{the-effective-nose-radius}

The radius fed to the correlation is not, in general, the nose cap
radius. For a sphere the only available length is its radius, so the
stagnation velocity gradient goes as 1/R and the 1/√R heating law
follows. For a shallow spherical cap the pressure falls from its
stagnation value over a distance set by the body radius and the corner,
so the benefit of flattening the nose saturates. The repair used by both
primaries is to keep the sphere formula and feed it the radius of the
hemisphere that would give the body's actual gradient {[}5, 6{]}. The
implementation interpolates Ellison's nine measured values of
R\_b/R\_eff, with an exact hemisphere anchor at K = 1, shape-preserving
cubic in K and linear in corner ratio. That interpolation is this
project's own construction (A-GEO-3a). For K \textgreater{} 1 the model
falls back to R\_eff = R\_n and flags the result. Validity from the
primaries: zero angle of attack, K in {[}0, 1{]}, R\_c/R\_b ≤ 0.4, Mach
above about 3.5.

The two models agree identically at K = 1. The burn-versus-bake results
of §7 use a hemisphere-nosed body (R\_n = D/2), so they are unchanged by
the correction, and a test reproduces the persisted results bit for bit
under the legacy default. The optimised designs of §10 sit at K ≈ 0.42,
where the correction raises peak flux by about 20\% (§17.1).

\subsection{Thermal protection system}

One-dimensional transient conduction through a multilayer stack: finite
volumes, harmonic-mean interface conductivities, backward Euler (A-TPS-1
to A-TPS-9). The harmonic mean is what series thermal resistances give,
and it matters most at the insulator-on-aluminium interface that sets
the bondline temperature. Backward Euler is unconditionally stable; the
explicit limit for this mesh is about 0.018 s, against roughly 2800
implicit steps per evaluation. It is first order in time, and
Crank--Nicolson remains an open item.

The surface balance, q\_net = q\_conv − εσ(T\_s\textsuperscript{4} −
T\_sink\textsuperscript{4}) with ε = 0.85 and a 4 K sink, is
Newton-linearised about the current surface estimate. The first
implementation used fixed-point iteration and diverged (NR-01). The back
face is adiabatic, which over-predicts the bondline for every case
alike. The solve continues for 1200 s past the end of the trajectory at
zero incident flux, because the bondline peak lags the pulse; omitting
that phase under-predicted the bondline peak by about 50 K on a 20 mm
stack (NR-02). Each candidate carries its bondline warming rate at the
end of the window so a truncated solve is visible.

Material properties are placeholders of the right order for a
low-density insulator on an aluminium-like structure (A-TPS-7). They are
not those of any qualified TPS material, so absolute temperatures are
indicative only. The insulator is 15 mm. At the original 40 mm the
bondline moved 8.8 K over a whole entry, because √(αt) with α ≈
8.9×10\textsuperscript{−7} m\textsuperscript{2}/s over a 245 s entry is
about 15 mm (NR-03). The effect studied here can be designed away with
insulation, at a mass cost.

\subsection{Geometry}

The forebody is parametrised as sphere, cone, torus, cone, flat base,
with a pure spherical segment as the limit of zero fore-cone length.
\texttt{CapsuleGeometry.validate()} decides validity inside the
evaluator. An invalid capsule is returned as an infeasible candidate
with the validator's message and stays in the log, so the boundary of
the valid region is data.

\subsection{Constraints and metrics}

Every evaluation returns peak heat flux, integrated external heat load,
peak surface temperature, peak bondline temperature, a bondline exposure
metric, thermal penetration depth, maximum deceleration, maximum dynamic
pressure, entry duration, feasibility and every constraint margin. An
unsourced limit is stored as \texttt{null} and skipped, never treated as
satisfied (A-LIM-2).

Three limits are active in the optimisation studies, and their standing
differs.

\begin{itemize}
\tightlist
\item
  Bondline allowable, 450 K. Sourced: the Shuttle's aluminium-structure
  design limit, 449.8 K {[}12{]}. It is a substrate limit, right here
  only because the modelled structure is aluminium-like (A-LIM-1a).
\item
  Deceleration, 12 g. Not sourceable as a flat number (A-LIM-1b). It
  matches no documented sustained-g curve. I have not changed it,
  because the choice between relabelling it as an emergency envelope and
  adopting the NASA-STD-3001 duration curve decides what kind of vehicle
  this is. It is an open decision (§16).
\item
  Heat-shield mass fraction, ≤ 1.0. A logical necessity, not a mass
  budget (A-OPT-6): the forebody TPS cannot outweigh the vehicle. It was
  added after the optimiser produced shields heavier than the vehicle
  (§17.1).
\end{itemize}

Each candidate also carries audit diagnostics: the share of heat load
accrued above 86 km, the soak-out warming rate, bluntness ratio and
heat-shield mass fraction.

\subsection{Fidelity 1 and the studies built on it}

The CFD model and its gate are §8; the drag surface is §9. Optimisation
methods and budgets are in §10, the optimiser comparison in §11, the
adaptive-fidelity design in §12 and the uncertainty design in §13. Each
study's success criterion was written into a configuration file before
the study ran, and each long run hashes the physics source at launch and
aborts if it changes (§17.4).

\section{Verification and Validation}

Verification asks whether the equations are solved correctly; validation
asks whether they are the right equations. Status vocabulary is fixed:
\texttt{NOT\_\allowbreak{}STARTED}, \texttt{IN\_\allowbreak{}PROGRESS},
\texttt{PASS}, \texttt{FAIL}, \texttt{LIMITED}. A row is \texttt{PASS}
only when checked against an independent reference, and \texttt{LIMITED}
is never written as \texttt{PASS}. The full matrix is
\texttt{VALIDATION\_\allowbreak{}MATRIX.md}; Table 1 summarises it.

\begin{table}[htbp]
\caption{Gate statuses (source: \texttt{VALIDATION\_\allowbreak{}MATRIX.md},
2026-09-21).}
\small\centering
\begin{tabular}{@{}>{\raggedright\arraybackslash}p{\dimexpr(\linewidth-10\tabcolsep)*61/1000\relax}>{\raggedright\arraybackslash}p{\dimexpr(\linewidth-10\tabcolsep)*134/1000\relax}>{\raggedright\arraybackslash}p{\dimexpr(\linewidth-10\tabcolsep)*276/1000\relax}>{\raggedright\arraybackslash}p{\dimexpr(\linewidth-10\tabcolsep)*276/1000\relax}>{\raggedright\arraybackslash}p{\dimexpr(\linewidth-10\tabcolsep)*251/1000\relax}@{}}
\toprule
\tabhead{Gate} & \tabhead{Subject} & \tabhead{Reference} & \tabhead{Measured} & \tabhead{Status} \\
\midrule
G1A & Atmosphere 0--86 km & published USSA-76 layer-base values {[}24{]} & within 0.01\% on T, 0.1\% on p & PASS \\
G1A{\symfont ′} & Atmosphere 86--150 km & primary Table I at 90, 100, 110, 120, 150 km {[}24{]} & agreement to every printed digit at those five & PASS at five altitudes, LIMITED elsewhere \\
G1B & Trajectory & Putnam \& Braun Table 2 {[}9{]} & nine figures within 0.92 points (allowance 1.0) & PASS \\
G2 & Heating constant & TR R-376 eq. (33), Table II {[}1{]} & re-derived to +0.39\% (primary's own error 3.3\%) & PASS \\
G2{\symfont ′} & Heating model form & none & not attempted & LIMITED \\
G3 & TPS conduction & semi-infinite solid, constant flux (see §20, second block) & 0.002\% at surface; energy residual \textasciitilde1e-14 & PASS \\
G3b & Bench boundary terms & closed forms & 1e-4 steady state, 2e-3 lumped decay & PASS (internal) \\
G-GEO & Capsule geometry & internal cross-checks & 42 tests & PASS (internal) / LIMITED \\
G-GEO2 & Effective nose radius & Ellison Table I {[}6{]}; Zoby \& Sullivan {[}5{]} & exact at nine points and both limits & PASS implementation / LIMITED physics \\
G4 & CFD & §8 & §8 & PASS, with a history (§8.2) \\
G5 & Coupled evaluator & self-equality of fingerprints & identical over three runs & PASS (software gate only) \\
M3 & Drag surface & none for capsules & §9 & LIMITED \\
M4--M7 & Studies & none external & §10--§13 & LIMITED \\
M8 & Coupon experiment & none yet & no measurement exists & IN\_PROGRESS \\
\bottomrule
\end{tabular}
\end{table}

\subsection{Atmosphere}

G1A{\symfont ′} was opened because a widely used online table labels its
data above 84.852 km as a different atmosphere model in a place that is
easy to miss. The transcribed values at five altitudes were diffed
against 300 dpi renders of the primary with tolerance set to half the
primary's printed resolution. They agree to every printed digit; the
largest relative difference is 1.6×10\textsuperscript{−15}. That does
not establish the 86, 95 and 130 km rows, the pressure column, or the
log-interpolation between rows. This matters because 5 to 20\% of a
feasible design's integrated heat load accrues above 86 km (§17.1), and
the bondline objective responds to integrated load.

\subsection{Trajectory}

The validation reproduces a published disagreement instead of seeking
agreement with the closed form. Under Putnam and Braun's stated
assumptions (exponential atmosphere with ρ\textsubscript{0} = 1.215
kg/m\textsuperscript{3} and 8500 m scale height, constant gravity 9.81
m/s\textsuperscript{2}, radius 6378 km, no lift) the integrator
reproduces all nine published figures. The largest residual is 0.92
percentage points, on sample-return peak deceleration (+35.82\% against
a published +34.9\%).

One artefact found on the way looked like physics. The steep case peaks
near 6 km altitude at 4.4 km/s, so consecutive output samples are about
440 m apart and the raw peak altitude was quantised; the first test
reported +5.16\% against a published +3.6\%. Parabolic refinement
through the three samples around the peak removed it. G1B validates the
integration under an exponential atmosphere, constant gravity and
constant C\_D. It does not validate the production configuration, and no
flight trajectory has been reproduced.

\subsection{Heating: the constant, and the model form}

The primary's equation and K = 0.1113 for air, carried through the three
substitutions with the units handled explicitly, give

\begin{center}
k = K\textsubscript{air} /\allowbreak{} (2√101325) = 0.1113 /\allowbreak{} 636.632 = 1.74826×10\textsuperscript{−4},
\end{center}

which is +0.39\% from the 1.74150×10\textsuperscript{−4} in the code.
The declared tolerance is the primary's own 3.3\%.

I did not change the constant, because the derivation is less accurate
than the gap it measures. At the baseline condition the exact
perfect-gas stagnation-pressure coefficient is 1.83937, not the
Newtonian 2, worth −4.10\% on k; the cold-wall assumption is worth
−1.10\%; the dropped freestream enthalpy +0.81\%. Applying all three
gives 1.6717×10\textsuperscript{−4}, which is −4.01\% in the opposite
direction. The code's value sits between two defensible derivations and
the evidence cannot separate them. The exact digits 1.7415 could not be
reconstructed from the primary by any route tried. The useful output is
a bound: treat the correlation as ±4\% at best and propagate that
(A-HEAT-1, NR-12).

G2{\symfont ′}, the model form, is \texttt{LIMITED} and was not
attempted: catalycity (roughly a factor of two {[}4{]}), the cold-wall
assumption applied to a surface near 2400 K, neglected radiation, and no
rejection of inputs outside the primary's fitted domain. Every
uncertainty figure in §13 is therefore a lower bound.

\subsection{Conduction}

G3 compares the solver with the analytical solution for a semi-infinite
solid under constant surface flux, at four depths and at the surface,
against a declared tolerance of 0.2\%. Measured: 0.002\% at the surface
and an energy residual near 1e-14, with grid and timestep convergence
demonstrated (\texttt{tests/\allowbreak{}test\_\allowbreak{}tps.py}). A
residual at that level is round-off, the expected signature of a
conservative finite-volume scheme in which each face flux appears once
with each sign. This is the \texttt{PASS} the headline result rests on.
The citation for the analytical solution has not itself been checked
against a copy of the book (§20).

\subsection{What the gates do not establish}

G1A{\symfont ′} validates a transcription at five altitudes. G1B
validates an integration under assumptions chosen so that a closed form
exists. G2 validates a constant, not a correlation. G4 validates a CFD
pipeline on a sphere at Mach 3 and 6 (§8). Nothing in the matrix is
validated against a physical measurement.

The sourcing pass that produced three of these rows changed no number.
The baseline was re-run before and after and is bit-identical: 160.28
W/cm\textsuperscript{2}, 494.8 K, 11.89 g. I had written beforehand that
at least one of the numbers would turn out wrong.

\section{Burn-vs-Bake Result}

Everything in this section is Fidelity 0: constant C\_D = 1.20, and the
legacy effective-radius setting R\_eff = R\_n.~For these geometries
(hemisphere nose, K = 1) the legacy and corrected heating models agree
identically, so the numbers are unaffected by the correction of §5.4.
They are affected by the constant drag coefficient: §9 shows that the
CFD-derived surface raises the baseline capsule's peak flux by 16.52\%
and its bondline by 9.56 K. Read the magnitudes here as legacy-model
numbers. Run \texttt{M1-\allowbreak{}20260902T130547Z}.

\subsection{The sweep}

Entry flight-path angle was swept from −8.00° to −1.50° in 41 steps with
every other variable at baseline. Each candidate ran the full chain;
none was discarded.

\begin{table}[htbp]
\caption{Endpoints of the sweep (source:
\texttt{reports/\allowbreak{}milestones/\allowbreak{}M1\_\allowbreak{}burn\_\allowbreak{}vs\_\allowbreak{}bake.md}).}
\small\centering
\begin{tabular}{@{}lll@{}}
\toprule
 & \tabhead{A: steep, γ\textsubscript{0} = −8.00°} & \tabhead{B: shallow, γ\textsubscript{0} = −1.50°} \\
\midrule
Peak heat flux & 230.64 W/cm\textsuperscript{2} & 142.38 W/cm\textsuperscript{2} \\
Integrated external load & 79.3 MJ/m\textsuperscript{2} & 128.3 MJ/m\textsuperscript{2} \\
Entry duration & 134.8 s & 324.1 s \\
Peak bondline temperature & 423.7 K & 537.8 K \\
\bottomrule
\end{tabular}
\end{table}

B reduces peak flux by 38.3\% and runs the bondline 114.1 K hotter. An
automated search over all 1640 ordered pairs, with the thresholds of §4,
found 818 counterexample pairs
(\texttt{M1\_\allowbreak{}signature\_\allowbreak{}counterexample.md}).

What this establishes is monotonicity. Over the whole range, every step
that lowers peak flux raises bondline temperature, and no interior angle
improves both. That is enough for H0. The Spearman rank correlation over
the sweep is −1.000, and that number is tautological: two strictly
monotone functions of one swept variable can only give ±1. It restates
``both are monotone in γ with opposite signs'' and measures nothing
about a design space, because this domain is a line segment.

\begin{figure}[htbp]
\centering
\includegraphics[width=\linewidth,height=0.42\textheight,keepaspectratio]{/Users/akj/Projects/aether/reports/final/../figures/M1_anticorrelation.pdf}
\caption{Peak stagnation heat flux (W/cm\textsuperscript{2}) and peak bondline
temperature (K) against entry flight-path angle (degrees), 41-point
sweep, Fidelity 0, run \texttt{M1-\allowbreak{}20260902T130547Z}. File
\texttt{reports/\allowbreak{}figures/\allowbreak{}M1\_\allowbreak{}anticorrelation.png}.}
\end{figure}

\subsection{Mechanism}

The shallow entry decelerates in thinner air for longer. Peak flux
falls, the pulse lengthens from 135 s to 324 s, and the integrated load
rises from 79.3 to 128.3 MJ/m\textsuperscript{2}. The stack's diffusion
length over a 245 s entry is about its own thickness. A short intense
pulse is absorbed near the surface and re-radiated at
T\textsuperscript{4} before it diffuses inward; a long mild pulse
reaches the bondline.

\begin{figure}[htbp]
\centering
\includegraphics[width=\linewidth,height=0.46\textheight,keepaspectratio]{/Users/akj/Projects/aether/reports/final/../figures/M1_mechanism.pdf}
\caption{Heat-flux histories and in-depth temperature response for the steep and
shallow entries of Table 2, Fidelity 0, run
\texttt{M1-\allowbreak{}20260902T130547Z}. File
\texttt{reports/\allowbreak{}figures/\allowbreak{}M1\_\allowbreak{}mechanism.png}.
The depth-against-time field is
\texttt{reports/\allowbreak{}figures/\allowbreak{}M1\_\allowbreak{}temperature\_\allowbreak{}field.png};
integrated load against bondline temperature is
\texttt{M1\_\allowbreak{}integrated\_\allowbreak{}vs\_\allowbreak{}bondline.png}.}
\end{figure}

\subsection{No candidate was feasible}

Zero of the 41 candidates satisfy both hard constraints. Steep entries
exceed 12 g; shallow ones exceed 450 K at the bondline. I did not
anticipate this. With geometry frozen, trajectory shape cannot produce a
valid design, which is why the geometry axis was opened immediately.

\subsection{Opening the diameter axis}

A full-factorial grid over entry angle and diameter, 140 evaluations
with mass and bluntness ratio held constant, recovers 21 feasible
designs.

\begin{table}[htbp]
\caption{Fidelity-0 optimiser comparison (source: M1 report §M1b).}
\small\centering
\begin{tabular}{@{}llll@{}}
\toprule
 & \tabhead{Baseline} & \tabhead{Peak-flux-only optimum} & \tabhead{Joint optimum} \\
\midrule
Entry angle & −3.00° & −1.50° & −3.50° \\
Diameter & 1.20 m & 3.00 m & 3.00 m \\
Ballistic coefficient & 257.9 kg/m\textsuperscript{2} & 41.3 & 41.3 \\
Peak heat flux & 160.28 W/cm\textsuperscript{2} & 37.01 & 44.27 \\
Integrated load & 112.5 MJ/m\textsuperscript{2} & 32.1 & 26.5 \\
Peak surface temperature & 2381 K & 1626 & 1699 \\
Peak bondline temperature & 494.8 K & 444.1 & 411.6 \\
Maximum deceleration & 11.89 g & 9.13 & 11.69 \\
Feasible & no & yes & yes \\
\bottomrule
\end{tabular}
\end{table}

The joint optimum runs 32.5 K cooler at the bondline for +19.6\% peak
flux. Both optima sit on the 3.0 m diameter bound, so the bound is
selecting them, and they should be read as ``at least this far in this
direction''. What happened when the bound was widened is in §17.1.

\begin{figure}[htbp]
\centering
\includegraphics[width=\linewidth,height=0.42\textheight,keepaspectratio]{/Users/akj/Projects/aether/reports/final/../figures/M1b_feasible_region.pdf}
\caption{Feasible region on the entry-angle (degrees) by diameter (m) grid with
the 12 g and 450 K constraints, 140 designs, Fidelity 0, run
\texttt{M1-\allowbreak{}20260902T130547Z}. File
\texttt{reports/\allowbreak{}figures/\allowbreak{}M1b\_\allowbreak{}feasible\_\allowbreak{}region.png}.
The two selected designs are compared in
\texttt{M1b\_\allowbreak{}optimiser\_\allowbreak{}comparison.png}; the
trade space of the one-parameter sweep is
\texttt{M1\_\allowbreak{}trade\_\allowbreak{}space.png}.}
\end{figure}

The deceleration limit decides this region. On the stored grid, a 10 g
limit leaves 4 of 140 designs feasible and an 8 g limit leaves none
(\texttt{ASSUMPTIONS.md}, open decision under A-LIM-1b).

\section{CFD}

\subsection{Model and scope}

\texttt{rhoCentralFoam}, axisymmetric Euler equations on a 5° wedge,
calorically perfect air (γ = 1.4), Kurganov flux with van Leer
reconstruction, local time stepping, forebody-only domain ending in a
supersonic outflow plane. OpenFOAM v2512 was used; the specification
names v2606 (A-CFD-6). The forebody-only domain was chosen because an
inviscid wake never reached a steady base pressure (NR-06), so base drag
is an assumption with a band (§9).

Within its tested assumptions the CFD may be used for forebody
pressure-drag coefficients, forebody surface pressure and perfect-gas
shock shape. It may not be used for heating of any kind, base drag,
flight shock stand-off, shock-layer temperatures, lift, moments or
stability
(\texttt{docs/\allowbreak{}validation/\allowbreak{}M2\_\allowbreak{}model\_\allowbreak{}form\_\allowbreak{}limits.md}).
No heating quantity in this paper comes from CFD.

\subsection{Gate G4}

Benchmark: a sphere of radius 0.5 m at Mach 3 and 6 on three meshes
(3,072, 12,288 and 49,152 cells). Run
\texttt{M2-\allowbreak{}20260920T123901Z}. The force criterion, declared
before the runs, requires forebody C\_D over the final 2000 iterations
to have peak-to-peak variation ≤ 0.1\% and half-window drift ≤ 0.02\%.

\begin{table}[htbp]
\caption{Mesh study and benchmark (source:
\texttt{reports/\allowbreak{}milestones/\allowbreak{}M2\_\allowbreak{}cfd\_\allowbreak{}validation.md}).}
\footnotesize\centering
\begin{tabular}{@{}>{\raggedright\arraybackslash}p{\dimexpr(\linewidth-14\tabcolsep)*50/1000\relax}>{\raggedright\arraybackslash}p{\dimexpr(\linewidth-14\tabcolsep)*169/1000\relax}>{\raggedright\arraybackslash}p{\dimexpr(\linewidth-14\tabcolsep)*88/1000\relax}>{\raggedright\arraybackslash}p{\dimexpr(\linewidth-14\tabcolsep)*69/1000\relax}>{\raggedright\arraybackslash}p{\dimexpr(\linewidth-14\tabcolsep)*169/1000\relax}>{\raggedright\arraybackslash}p{\dimexpr(\linewidth-14\tabcolsep)*220/1000\relax}>{\raggedright\arraybackslash}p{\dimexpr(\linewidth-14\tabcolsep)*232/1000\relax}@{}}
\toprule
\tabhead{Mach} & \tabhead{C\_D coarse / medium / fine} & \tabhead{observed order} & \tabhead{GCI\_fine} & \tabhead{p\textsubscript{0}/p∞ vs Rayleigh pitot} & \tabhead{stand-off vs Van Dyke \& Gordon {[}26{]}} & \tabhead{C\_D,fore vs Clark's expression {[}27{]}} \\
\midrule
3 & 0.84440 / 0.84613 / 0.84692 & 1.13 & 0.098\% & +0.14\% (±1\%) & +0.95\% (±5\%) & −0.32\% (±3\%) \\
6 & 0.86890 / 0.87090 / 0.87199 & 0.88 & 0.187\% & −0.44\% (±1\%) & +0.65\% (±5\%) & −1.82\% (±3\%) \\
\bottomrule
\end{tabular}
\end{table}

The gate reads \texttt{PASS}, and it has a history that belongs next to
it. On the original runs both fine-mesh cases ended in a bounded limit
cycle, peak-to-peak 0.282\% and 0.223\% after 100,000 iterations with
drift one to two orders below its limit, and the gate read
\texttt{LIMITED}
(\texttt{gate\_\allowbreak{}assessment\_\allowbreak{}20260921\_\allowbreak{}0057\_\allowbreak{}before\_\allowbreak{}restarts.json},
kept). The criterion was not changed. Following the specification's
failure-handling order, each fine case was continued as a new case at a
lower Courant limit (0.1), and met the criterion in two consecutive
blocks (0.033\% and 0.077\%). The restart rule and the rule for which
solution is ``of record'' were written after the original results had
been seen (NR-25). A reader who does not accept a rule written after the
fact should read G4 as \texttt{LIMITED}; the evidence for both readings
is in the run directory.

Other things a reader should know. The observed order at Mach 6 is below
one, and it is computed from a fine-to-medium difference of about 0.1\%,
so it is weakly determined. Stagnation pressure and stand-off are single
final-iteration snapshots; for p\textsubscript{0}/p∞ the snapshot noise
is as large as the mesh-to-mesh differences, so its Richardson value is
not relied on. The Mach 6 convergence margin is thin. Sphere total drag
is not validated: the measured free-flight totals are forebody plus base
plus friction, and this CFD computes the first only. A full-body
negative case was run inside the pipeline and fails the criterion, as it
should.

\begin{figure}[htbp]
\centering
\includegraphics[width=\linewidth,height=0.70\textheight,keepaspectratio]{/Users/akj/Projects/aether/reports/final/../figures/cfd/M2_sphere_M6_fine_fourpanel.pdf}
\caption{Sphere, R = 0.5 m, M∞ = 6, fine mesh (refinement factor 4, 49,152
cells), the max-Courant-0.1 solution of record after NR-25, saved
iteration 110,000: Mach number, p/p∞, T/T∞ and ρ/ρ∞ in the meridional
plane, one flat-shaded polygon per cell, run
\texttt{M2-\allowbreak{}20260920T123901Z}, case
\texttt{sphere\_\allowbreak{}M6\_\allowbreak{}fine\_\allowbreak{}Co0p1}.
rhoCentralFoam, axisymmetric Euler, calorically perfect gas γ = 1.4, p∞
= 1000 Pa, T∞ = 220 K. White line: sonic line M = 1 from the cell-centre
triangulation; it bounds the subsonic nose region and also runs along
the captured shock, where M passes through 1 inside the one to two cells
of numerical shock thickness. Dashed: Billig's correlation for the shock
shape, an independent empirical curve, not a fit. Δ: shock stand-off
from \texttt{case\_\allowbreak{}result.json} (50\% density-rise point on
the stagnation line). Colour bars: viridis (Mach, ρ/ρ∞), magma (p/p∞),
cividis (T/T∞). Within its tested assumptions this CFD is used for
forebody pressure drag, surface pressure and perfect-gas shock shape
only; it is not used for shock-layer temperatures or heating of any kind
(§8.1). File
\texttt{reports/\allowbreak{}figures/\allowbreak{}cfd/\allowbreak{}M2\_\allowbreak{}sphere\_\allowbreak{}M6\_\allowbreak{}fine\_\allowbreak{}fourpanel.png}.}
\end{figure}

\begin{figure}[htbp]
\centering
\includegraphics[width=\linewidth,height=0.42\textheight,keepaspectratio]{/Users/akj/Projects/aether/reports/final/../figures/cfd/M2_sphere_M6_mesh_levels.pdf}
\caption{Mach number for the sphere at M∞ = 6 on the coarse, medium and fine
meshes (refinement factors 1, 2, 4; saved iterations 10,000, 20,000 and
110,000), common colour scale (viridis, 0 to M∞), run
\texttt{M2-\allowbreak{}20260920T123901Z}, cases
\texttt{sphere\_\allowbreak{}M6\_\allowbreak{}coarse},
\texttt{sphere\_\allowbreak{}M6\_\allowbreak{}medium},
\texttt{sphere\_\allowbreak{}M6\_\allowbreak{}fine\_\allowbreak{}Co0p1}.
Cell edges are drawn on the coarse and medium meshes so the refinement
is visible; on the fine mesh (49,152 cells) they would blacken the panel
and are omitted. The fine panel is the max-Courant-0.1 restart of record
(NR-25); coarse and medium ran at max Courant 0.2. C\_D,fore and Δ/R are
read from each \texttt{case\_\allowbreak{}result.json}. Same gas model
and freestream as Figure 6. File
\texttt{reports/\allowbreak{}figures/\allowbreak{}cfd/\allowbreak{}M2\_\allowbreak{}sphere\_\allowbreak{}M6\_\allowbreak{}mesh\_\allowbreak{}levels.png}.}
\end{figure}

\begin{figure}[htbp]
\centering
\includegraphics[width=\linewidth,height=0.46\textheight,keepaspectratio]{/Users/akj/Projects/aether/reports/final/../figures/cfd/M2_sphere_stagnation_line_profiles.pdf}
\caption{Sampled cell values of p/p∞, T/T∞, ρ/ρ∞ and Mach number along the
symmetry axis (\texttt{postProcess sampleLine}) for the sphere at M∞ = 3
and 6 on the three meshes, each at its final iteration (15,000, 20,000
and 110,000 at Mach 3; 10,000, 20,000 and 110,000 at Mach 6; the fine
curves are the max-Courant-0.1 restarts of record), run
\texttt{M2-\allowbreak{}20260920T123901Z}. Thin vertical lines: the
shock stand-off Δ of each mesh (50\% density-rise point,
\texttt{case\_\allowbreak{}result.json}); dashed: the Billig (1967)
stand-off correlation, which is cited but not verified in primary by
this project (§20). Line style and colour distinguish the meshes. Same
gas model and freestream as Figure 6. File
\texttt{reports/\allowbreak{}figures/\allowbreak{}cfd/\allowbreak{}M2\_\allowbreak{}sphere\_\allowbreak{}stagnation\_\allowbreak{}line\_\allowbreak{}profiles.png}.}
\end{figure}

\begin{figure}[htbp]
\centering
\includegraphics[width=\linewidth,height=0.46\textheight,keepaspectratio]{/Users/akj/Projects/aether/reports/final/../figures/M2_mesh_convergence.pdf}
\caption{Forebody drag coefficient, stand-off distance and stagnation pressure
ratio against mesh level for the sphere at Mach 3 and 6, run
\texttt{M2-\allowbreak{}20260920T123901Z}. File
\texttt{reports/\allowbreak{}figures/\allowbreak{}M2\_\allowbreak{}mesh\_\allowbreak{}convergence.png}.
Force histories against the declared criterion:
\texttt{M2\_\allowbreak{}force\_\allowbreak{}convergence.png}; the
fine-mesh limit cycle:
\texttt{M2\_\allowbreak{}limit\_\allowbreak{}cycle.png}; benchmark
comparison: \texttt{M2\_\allowbreak{}benchmark.png}; residuals:
\texttt{M2\_\allowbreak{}residuals.png}.}
\end{figure}

The GCI follows the three-grid procedure with safety factor 1.25, with
formulae as given in NASA's NPARC verification tutorial {[}25{]}. The
paper usually cited for the procedure could not be opened (§20).

\section{Coupled Model}

\subsection{The drag surface}

C\_D(M, shape) = C\_D,fore(M, shape) + C\_D,base(M). The first term is a
Gaussian process (Matérn-5/2, one length scale per input) through
converged CFD cases over log Mach and three forebody ratios: bluntness
R\_n/D, cone half-angle, shoulder ratio R\_c/D. The second is a stated
assumption. Angle of attack is zero by construction. The trajectory
consumes a per-capsule one-dimensional interpolant and never calls CFD
or the GP inside a time step. Surface
\texttt{cfd\_\allowbreak{}surface\_\allowbreak{}v2}, hash
\texttt{b63b68a40321f38b}, run
\texttt{M3-\allowbreak{}DP-\allowbreak{}20260920T1610Z}.

Of 76 planned design points, 69 were usable, 5 were rejected for missing
the unchanged force criterion and 2 crashed; all stay on disk. Every
point is on the coarse mesh, not the medium mesh first intended, because
a medium case cost about eight times a coarse one on the available
laptop. Six coarse-to-medium pairs on capsule shapes differ by at most
0.411\% in C\_D,fore, which shows the size of the effect and does not
bound the error. The discretisation band, 0.628\%, is transferred from
the sphere's GCI by assumption (A-CFD-9). Scale invariance was measured:
the same shape at 1.2 m and 3 m gives identical C\_D,fore to the digits
printed.

\begin{figure}[htbp]
\centering
\includegraphics[width=\linewidth,height=0.70\textheight,keepaspectratio]{/Users/akj/Projects/aether/reports/final/../figures/cfd/M3_design_point_grid_mach.pdf}
\caption{Mach-number fields for six design-point shapes at M∞ = 3, 6, 20 and 27
on the coarse mesh (refinement factor 1, 3,072 cells), forebody domain
from the nose to the maximum-radius station closed by a supersonic
outflow, D = 1.2 m, each panel at its case's final saved iteration (in
its title), run \texttt{M3-\allowbreak{}DP-\allowbreak{}20260920T1610Z}.
Colour: viridis, one scale per column (0 to M∞); thin white line: sonic
line; C\_D,fore from each \texttt{case\_\allowbreak{}result.json}. Rows
5 and 6 are the two shapes that stayed \texttt{REJECTED} under NR-27 at
Mach 20 (force criterion never met in two consecutive blocks at Courant
0.2, 0.1 and 0.05) and, as NR-28 records, at Mach 27; their fields are
the final saved solution of the last Courant restart, shown for what
they are and not on the drag surface. There is no Mach-12 column: the
design sampled Mach 12 (and Mach 6, except the baseline scale-check case
\texttt{sc1p2}) only at fill-point shapes, so no usable case exists for
these six shapes there. \texttt{dp062} at Mach 27 crashed in all three
domain attempts without writing a second-order field (NR-28; Figure
A.4). Outlines are drawn to the maximum-radius station only, where the
domain ends. Same gas model and freestream as Figure 6; a calorically
perfect γ = 1.4 gas is the wrong gas at Mach 20 to 27 and is carried as
a declared ±5\% band on C\_D,fore (A-CFD-12). File
\texttt{reports/\allowbreak{}figures/\allowbreak{}cfd/\allowbreak{}M3\_\allowbreak{}design\_\allowbreak{}point\_\allowbreak{}grid\_\allowbreak{}mach.png}.}
\end{figure}

\begin{table}[htbp]
\caption{Surface accuracy (source:
\texttt{reports/\allowbreak{}milestones/\allowbreak{}M3\_\allowbreak{}coupled\_\allowbreak{}model.md}
§6; C\_D,fore spans 0.352 to 1.553).}
\small\centering
\begin{tabular}{@{}>{\raggedright\arraybackslash}p{\dimexpr(\linewidth-12\tabcolsep)*206/1000\relax}>{\raggedright\arraybackslash}p{\dimexpr(\linewidth-12\tabcolsep)*126/1000\relax}>{\raggedright\arraybackslash}p{\dimexpr(\linewidth-12\tabcolsep)*142/1000\relax}>{\raggedright\arraybackslash}p{\dimexpr(\linewidth-12\tabcolsep)*142/1000\relax}>{\raggedright\arraybackslash}p{\dimexpr(\linewidth-12\tabcolsep)*190/1000\relax}>{\raggedright\arraybackslash}p{\dimexpr(\linewidth-12\tabcolsep)*190/1000\relax}@{}}
\toprule
\tabhead{Set} & \tabhead{n} & \tabhead{RMSE} & \tabhead{R\textsuperscript{2}} & \tabhead{z-score s.d.} & \tabhead{68\% coverage} \\
\midrule
Held-out test & 8 & 0.0234 & 0.9958 & 1.70 & 0.25 \\
k-fold, all & 69 & 0.0130 & 0.9983 & 1.13 & 0.74 \\
\bottomrule
\end{tabular}
\end{table}

The GP is overconfident. A sigma inflation factor of 1.21 is stored with
the surface and applied before sampling; the held-out spread of 1.70 on
n = 8 is larger than that. A shape outside the convex hull of the CFD
points is refused and returned as a rejected candidate, so failed CFD
cases shrink the usable design space.

The baseline entry reaches Mach 27.1, and 55\% of its heat load accrues
above Mach 20, so the design includes a Mach 27 node (9 of 14 anchors
usable, NR-28). Running a perfect gas to Mach 27 closes an extrapolation
in Mach number. It does nothing about model form: a calorically perfect
γ = 1.4 gas is the wrong gas at these speeds, and that is carried as a
declared ±5\% half-band on C\_D,fore which this project has not
validated (A-CFD-12). Base drag is bounded, not predicted, from the
vacuum limit and a base-pressure ratio read off a graph in Miller
{[}20{]}.

\subsection{What a constant drag coefficient was hiding}

\begin{table}[htbp]
\caption{Same designs, constant C\_D against the surface (source: M3 report §10).}
\small\centering
\begin{tabular}{@{}>{\raggedright\arraybackslash}p{\dimexpr(\linewidth-12\tabcolsep)*230/1000\relax}>{\raggedright\arraybackslash}p{\dimexpr(\linewidth-12\tabcolsep)*95/1000\relax}>{\raggedright\arraybackslash}p{\dimexpr(\linewidth-12\tabcolsep)*214/1000\relax}>{\raggedright\arraybackslash}p{\dimexpr(\linewidth-12\tabcolsep)*150/1000\relax}>{\raggedright\arraybackslash}p{\dimexpr(\linewidth-12\tabcolsep)*158/1000\relax}>{\raggedright\arraybackslash}p{\dimexpr(\linewidth-12\tabcolsep)*150/1000\relax}@{}}
\toprule
\tabhead{Design} & \tabhead{C\_D constant} & \tabhead{C\_D surface at peak heating} & \tabhead{peak flux} & \tabhead{bondline} & \tabhead{heat load} \\
\midrule
Baseline capsule & 1.200 & 0.8899 & +16.52\% & +9.56 K & +16.08\% \\
Five Fidelity-0 front designs & 1.200 & 1.3634--1.3643 & −6.02 to −5.89\% & −3.15 to −2.46 K & −6.37 to −6.19\% \\
\bottomrule
\end{tabular}
\end{table}

One placeholder number was wrong in two directions at once, and the
direction depended on shape. The substitution penalises the
slender-to-moderate baseline, whose drag had been over-estimated, and
mildly rewards blunt shapes.

Over 12 swept shapes inside the hull at fixed diameter, mass and entry
state, C\_D at peak heating ranges from 0.496 to 1.378, which acting
through drag alone moves peak flux by −6.35\% to +60.84\% and bondline
temperature by −4.2 to +29.7 K. The model predicts that shape now
matters. How much of that survives a finer mesh and real-gas effects is
not established here.

\begin{figure}[htbp]
\centering
\includegraphics[width=\linewidth,height=0.42\textheight,keepaspectratio]{/Users/akj/Projects/aether/reports/final/../figures/M3_constant_vs_surface.pdf}
\caption{Change in peak heat flux (\%), bondline temperature (K) and maximum
deceleration (\%) when the constant C\_D = 1.20 is replaced by
\texttt{cfd\_\allowbreak{}surface\_\allowbreak{}v2}, for the baseline
capsule and five Fidelity-0 front designs, run
\texttt{M3-\allowbreak{}DP-\allowbreak{}20260920T1610Z}. File
\texttt{reports/\allowbreak{}figures/\allowbreak{}M3\_\allowbreak{}constant\_\allowbreak{}vs\_\allowbreak{}surface.png}.
Design coverage:
\texttt{M3\_\allowbreak{}design\_\allowbreak{}coverage.png}; C\_D
against Mach:
\texttt{M3\_\allowbreak{}cd\_\allowbreak{}vs\_\allowbreak{}mach.png};
cross-validation:
\texttt{M3\_\allowbreak{}surface\_\allowbreak{}cv.png}; shape sweep:
\texttt{M3\_\allowbreak{}shape\_\allowbreak{}sweep.png}; base-drag
share:
\texttt{M3\_\allowbreak{}base\_\allowbreak{}drag\_\allowbreak{}fraction.png};
capsule family:
\texttt{M3\_\allowbreak{}capsule\_\allowbreak{}family.png}.}
\end{figure}

Gate G5 is a software gate. It says the coupled evaluator is
deterministic and labels its aerodynamics. It says nothing about whether
the drag is right.

\section{Optimization}

Fidelity 1. DOE run
\texttt{M4-\allowbreak{}DOE-\allowbreak{}20260920T204844Z}, optimisation
run \texttt{M4-\allowbreak{}OPT-\allowbreak{}20260920T205252Z}. An
earlier Fidelity-0 version of this study is archived unchanged
(\texttt{M4\_\allowbreak{}pareto\_\allowbreak{}optimisation\_\allowbreak{}fidelity0.md});
its fronts were computed under the legacy heating model behind a
bluntness cap since removed, and nothing here reuses them.

\subsection{Design space and screening}

Seven design variables (diameter, bluntness ratio, shoulder ratio, cone
half-angle, aft cone angle, fineness ratio, entry angle), two mission
givens (entry velocity, mass) and one deferred variable (insulator
thickness, which with no mass objective would simply park on its upper
bound). One coupled evaluation takes a median 0.150 s, so variance-based
sensitivity was affordable. Of 3000 Latin Hypercube samples, 1328 were
geometrically valid and 186 feasible.

Sobol' indices (Saltelli design, 6144 evaluations, output centred;
estimator citations in §20) were computed on a sub-box in which every
capsule exists, covering 11.8\% of the shape box, and they describe only
that fraction. Peak flux is dominated by diameter (S\_T 0.896); bondline
temperature by diameter (0.429), insulator thickness (0.377) and entry
angle (0.239); deceleration by entry angle (0.979). The freeze rule,
declared before the run, freezes a design variable whose total-order
upper confidence bound is below 0.01 on every output and whose validity
index is below 0.03. It left four active variables: diameter, bluntness,
cone half-angle (which moves no objective but decides validity) and
entry angle. It froze the shoulder ratio, a decision §10.3 returns to.

\begin{figure}[htbp]
\centering
\includegraphics[width=\linewidth,height=0.42\textheight,keepaspectratio]{/Users/akj/Projects/aether/reports/final/../figures/M4_doe_sobol.pdf}
\caption{Sobol' first- and total-order indices with 95\% bootstrap intervals for
peak heat flux, peak bondline temperature, maximum deceleration and
heat-shield mass fraction, Saltelli sub-box, Fidelity 1, run
\texttt{M4-\allowbreak{}DOE-\allowbreak{}20260920T204844Z}. File
\texttt{reports/\allowbreak{}figures/\allowbreak{}M4\_\allowbreak{}doe\_\allowbreak{}sobol.png}.
One-at-a-time sweeps:
\texttt{M4\_\allowbreak{}doe\_\allowbreak{}oat.png}; Latin Hypercube
cloud: \texttt{M4\_\allowbreak{}doe\_\allowbreak{}lhs.png}.}
\end{figure}

\subsection{Methods and fronts}

Two objectives, both minimised: peak heat flux and peak bondline
temperature. The TPI was studied and discarded (§17.2), so integrated
load is reported on every design and not optimised. Three methods each
received exactly 1000 evaluations per seed over 7 seeds, through one
budget meter in which an invalid geometry costs one evaluation and a
repeat costs none.

\begin{table}[htbp]
\caption{Method comparison, normalised hypervolume against the fixed reference
point (2.5×10\textsuperscript{6} W/m\textsuperscript{2}, 450 K) (source:
M4 report §8).}
\small\centering
\begin{tabular}{@{}>{\raggedright\arraybackslash}p{\dimexpr(\linewidth-8\tabcolsep)*323/1000\relax}>{\raggedright\arraybackslash}p{\dimexpr(\linewidth-8\tabcolsep)*303/1000\relax}>{\raggedright\arraybackslash}p{\dimexpr(\linewidth-8\tabcolsep)*205/1000\relax}>{\raggedright\arraybackslash}p{\dimexpr(\linewidth-8\tabcolsep)*166/1000\relax}@{}}
\toprule
\tabhead{Method} & \tabhead{Final hypervolume, mean ± s.d.} & \tabhead{Feasible found (mean)} & \tabhead{Front size (mean)} \\
\midrule
Latin Hypercube search & 0.2874 ± 0.0214 & 38 & 3.0 \\
NSGA-II & 0.3467 ± 0.0055 & 703 & 32.1 \\
Scalarised differential evolution & 0.2997 ± 0.0150 & 112 & 4.3 \\
\bottomrule
\end{tabular}
\end{table}

No significance test was applied; with seven seeds, overlapping ranges
mean no measured difference. The scalarised baseline's number is a
statement about its declared configuration at this budget. The combined
front has 78 designs.

\begin{table}[htbp]
\caption{Selected designs from the combined feasible front (source: M4 report
§10).}
\footnotesize\centering
\begin{tabular}{@{}>{\raggedright\arraybackslash}p{\dimexpr(\linewidth-18\tabcolsep)*160/1000\relax}>{\raggedright\arraybackslash}p{\dimexpr(\linewidth-18\tabcolsep)*124/1000\relax}>{\raggedright\arraybackslash}p{\dimexpr(\linewidth-18\tabcolsep)*124/1000\relax}>{\raggedright\arraybackslash}p{\dimexpr(\linewidth-18\tabcolsep)*131/1000\relax}>{\raggedright\arraybackslash}p{\dimexpr(\linewidth-18\tabcolsep)*138/1000\relax}>{\raggedright\arraybackslash}p{\dimexpr(\linewidth-18\tabcolsep)*58/1000\relax}>{\raggedright\arraybackslash}p{\dimexpr(\linewidth-18\tabcolsep)*94/1000\relax}>{\raggedright\arraybackslash}p{\dimexpr(\linewidth-18\tabcolsep)*87/1000\relax}>{\raggedright\arraybackslash}p{\dimexpr(\linewidth-18\tabcolsep)*80/1000\relax}@{}}
\toprule
\tabhead{Design} & \tabhead{Peak flux {[}W/m\textsuperscript{2}{]}} & \tabhead{Heat load {[}J/m\textsuperscript{2}{]}} & \tabhead{Peak surface T {[}K{]}} & \tabhead{Peak bondline T {[}K{]}} & \tabhead{Max g} & \tabhead{Entry angle} & \tabhead{Diameter {[}m{]}} & \tabhead{Feasible} \\
\midrule
Reference, unoptimised & 1.868×10\textsuperscript{6} & 1.306×10\textsuperscript{8} & 2476.1 & 504.3 & 12.25 & −3.00° & 1.2 & no \\
Peak-flux-only optimum & 2.295×10\textsuperscript{5} & 1.975×10\textsuperscript{7} & 1427.7 & 418.4 & 9.22 & −1.502° & 3.37 & yes \\
Joint knee & 2.510×10\textsuperscript{5} & 1.781×10\textsuperscript{7} & 1458.7 & 403.5 & 10.35 & −2.497° & 3.367 & yes \\
Bondline-only optimum & 2.783×10\textsuperscript{5} & 1.621×10\textsuperscript{7} & 1495.7 & 391.7 & 11.98 & −3.565° & 3.355 & yes \\
\bottomrule
\end{tabular}
\end{table}

Relative to the peak-flux-only optimum, the knee's bondline runs 14.9 K
cooler and its peak flux is 2.146×10\textsuperscript{4}
W/m\textsuperscript{2} higher. The whole front spans 26.6 K and
4.882×10\textsuperscript{4} W/m\textsuperscript{2}. The model predicts a
trade-off, so H1 is supported within the tested assumptions: a search
that looked at peak flux alone would have stopped at the
hottest-bondline end of the front.

\begin{figure}[htbp]
\centering
\includegraphics[width=\linewidth,height=0.42\textheight,keepaspectratio]{/Users/akj/Projects/aether/reports/final/../figures/M4_pareto_front.pdf}
\caption{Combined feasible Pareto front, peak heat flux (W/m\textsuperscript{2})
against peak bondline temperature (K), with the peak-flux-only optimum,
the knee and the bondline-only optimum marked; Fidelity 1
(\texttt{cfd\_\allowbreak{}surface\_\allowbreak{}v2}), run
\texttt{M4-\allowbreak{}OPT-\allowbreak{}20260920T205252Z}. File
\texttt{reports/\allowbreak{}figures/\allowbreak{}M4\_\allowbreak{}pareto\_\allowbreak{}front.png}.
Hypervolume histories: \texttt{M4\_\allowbreak{}hypervolume.png}; front
in design space:
\texttt{M4\_\allowbreak{}front\_\allowbreak{}variables.png}.}
\end{figure}

\subsection{The audit: what is selecting these designs}

The study includes a metric-gaming audit on the 78 front designs. Its
findings decide how Table 8 should be read.

Every front diameter lies between 3.355 and 3.376 m, and 50 of 78
designs are within 1\% of the heat-shield mass-fraction limit. The knee
is a 3.37 m, 350 kg vehicle whose forebody TPS weighs 346 kg. That is
not a capsule. Vehicle mass is fixed while diameter varies, so the
optimiser buys a low ballistic coefficient (about 29
kg/m\textsuperscript{2}) with a shield no real vehicle could carry,
until the one constraint that knows about mass stops it. There is no
sourced mass budget in this project. Removing the constraint gives a
front of six designs, all of which violate it.

28 of 78 designs sit on the edge of the CFD hull in bluntness, at R\_n/D
near 1.2. The hull ends there because that is where CFD anchors were
placed, and they were placed there because the Fidelity-0 front sat on
an earlier bluntness cap at 1.2. Labelled GP extrapolations beyond the
hull are worth about 1\% of peak flux, because the corrected nose model
has nearly saturated.

64 of 78 designs have the cone half-angle within 1\% of its 70° bound,
held in a window of about a degree between geometric validity and the
box.

The shoulder ratio was frozen at 0.10 by the screening rule, on a
total-order upper bound of 0.0048. Probes through the canonical
evaluator, entering no front, show that moving the three selected
designs to the box minimum of 0.02 lowers peak flux by 16.6 to 19.9
kW/m\textsuperscript{2}, about 7\%, and bondline temperature by 3.1 to
3.9 K (NR-29). A Sobol' index is a share of variance over the sub-box,
where diameter carries about 0.90 of the peak-flux variance, so a lever
worth 7\% at the front rounds to zero. The lever is also one this
heating model rewards for the wrong reason: it sees the larger effective
radius a sharp corner gives and cannot see the corner's own heating. The
freeze stands, labelled as a fence that happens to be in the right
place.

Every front design accrues more than 5\% of its heat load above 86 km,
up to 20.9\%.

Surface C\_D at peak heating spans 1.363 to 1.364 across the front,
about 0.1\%. Capsule shape, which Fidelity 1 was built to let matter,
does not vary along this front.

\subsection{What H1 amounts to}

Diameter is set by the mass fence, bluntness by the hull edge, cone
angle by validity and the box, shoulder by the freeze. The only variable
that trades the two objectives along the front is entry flight-path
angle, from −1.50° (a box bound) to −3.57° (the 12 g limit). The
Fidelity-1 front is a burn-versus-bake curve in entry angle drawn at a
geometry chosen by fences. H1 is supported within the model by the same
mechanism as §7. It is not evidence that joint optimisation finds a
better capsule shape.

\begin{figure}[htbp]
\centering
\includegraphics[width=\linewidth,height=0.42\textheight,keepaspectratio]{/Users/akj/Projects/aether/reports/final/../figures/infographics/03_pareto_front_fences.pdf}
\caption{What the optimiser found: the combined feasible front of §10.2 is an
entry-angle curve at a geometry fixed by fences, with the robust knee of
§13.5 and its 95th-percentile whiskers. Explanatory figure. The fences
are drawn in objective space as pickets hugging the front; they are
constraints in design space and the picket positions are illustrative,
while the counts on them are the audit's (§10.3). Sources:
\texttt{results/\allowbreak{}M4/\allowbreak{}M4-\allowbreak{}OPT-\allowbreak{}20260920T205252Z/\allowbreak{}candidates.csv}
and \texttt{summary.json},
\texttt{results/\allowbreak{}M7/\allowbreak{}M7-\allowbreak{}LFL-\allowbreak{}20260921T011854Z/\allowbreak{}likeforlike.json},
\texttt{results/\allowbreak{}M7/\allowbreak{}M7-\allowbreak{}UQ-\allowbreak{}20260920T233048Z/\allowbreak{}summary.json},
\texttt{results/\allowbreak{}M7/\allowbreak{}M7-\allowbreak{}ROBUST-\allowbreak{}20260920T211216Z/\allowbreak{}robust\_\allowbreak{}front.csv},
\texttt{reports/\allowbreak{}milestones/\allowbreak{}M4\_\allowbreak{}pareto\_\allowbreak{}optimisation.md}.
File
\texttt{reports/\allowbreak{}figures/\allowbreak{}infographics/\allowbreak{}03\_\allowbreak{}pareto\_\allowbreak{}front\_\allowbreak{}fences.png}.}
\end{figure}

\section{AI/Surrogate Comparison}

Fidelity 1. Run
\texttt{M5-\allowbreak{}ABL-\allowbreak{}20260920T211134Z}. The study
compared a large-language-model engineering agent with conventional
optimisers at 200 evaluations per seed over 5 seeds, on the four active
variables of §10. No CFD solver ran during this study; evaluations went
through the drag surface.

\subsection{Rule, methods, asymmetries}

Declared before any run
(\texttt{configs/\allowbreak{}ai\_\allowbreak{}ablation.yaml}): at 50,
100 and 200 evaluations the agent is compared with the non-LLM method of
highest mean hypervolume. ``AI helped'' requires a mean gain of at least
0.005 and a one-sided exact permutation test at p \textless{} 0.05 after
Holm correction; ``AI hurt'' is the mirror; anything else is ``no
measured difference''.

Methods: Latin Hypercube, NSGA-II at two populations, a Gaussian-process
Bayesian optimiser with ParEGO scalarisation
(\texttt{bo\_\allowbreak{}parego}), and the agent (model
\texttt{claude-\allowbreak{}sonnet-\allowbreak{}5}), which proposes
batches from structured tables under a strict schema. The agent and
\texttt{bo\_\allowbreak{}parego} start from the same seeded initial
design. The agent ran in an empty directory with tools and file
discovery disabled and was told neither the model's equations nor any
earlier finding.

Three asymmetries are part of what is measured. Variable and metric
names carry aerospace meaning, so the agent has prior knowledge the
numeric optimisers lack. The agent reads the geometry validator's
failure messages, which name the variable to move; the numeric
optimisers see only a violation. And the model may have seen similar
problems in training.

Two disclosures. The Bayesian optimiser was changed twice after a
one-seed smoke test and before any study run, and that seed was removed
(NR-17). An earlier study run was aborted when the physics changed under
it, and no number from it was inspected (NR-18).

\subsection{Result}

\begin{table}[htbp]
\caption{Pre-declared comparison (source:
\texttt{reports/\allowbreak{}milestones/\allowbreak{}M5\_\allowbreak{}ai\_\allowbreak{}ablation.md}
§4).}
\small\centering
\begin{tabular}{@{}>{\raggedright\arraybackslash}p{\dimexpr(\linewidth-12\tabcolsep)*142/1000\relax}>{\raggedright\arraybackslash}p{\dimexpr(\linewidth-12\tabcolsep)*275/1000\relax}>{\raggedright\arraybackslash}p{\dimexpr(\linewidth-12\tabcolsep)*132/1000\relax}>{\raggedright\arraybackslash}p{\dimexpr(\linewidth-12\tabcolsep)*132/1000\relax}>{\raggedright\arraybackslash}p{\dimexpr(\linewidth-12\tabcolsep)*91/1000\relax}>{\raggedright\arraybackslash}p{\dimexpr(\linewidth-12\tabcolsep)*224/1000\relax}@{}}
\toprule
\tabhead{Evaluations} & \tabhead{Best conventional (mean HV)} & \tabhead{Agent mean HV} & \tabhead{Difference} & \tabhead{Holm p} & \tabhead{Verdict} \\
\midrule
50 & \texttt{bo\_\allowbreak{}parego} 0.3145 & 0.3428 & +0.0283 & 0.0278 & AI helped \\
100 & \texttt{bo\_\allowbreak{}parego} 0.3465 & 0.3550 & +0.0085 & 0.0119 & AI helped \\
200 & \texttt{bo\_\allowbreak{}parego} 0.3520 & 0.3556 & +0.0036 & 0.0119 & no measured difference \\
\bottomrule
\end{tabular}
\end{table}

At 200 evaluations the test rejects (every agent seed is above every
\texttt{bo\_\allowbreak{}parego} seed), and the verdict is ``no measured
difference'' because the mean gain is under the threshold declared in
advance. The agent was distinguishably ahead by an amount I had declared
too small to count. With n = 5 the smallest attainable one-sided p is
1/252, so only complete separation can reject. NSGA-II at 200
evaluations reached 0.2803 ± 0.0253; the same algorithm reached 0.3467
after 1000 evaluations in §10.

\begin{figure}[htbp]
\centering
\includegraphics[width=\linewidth,height=0.42\textheight,keepaspectratio]{/Users/akj/Projects/aether/reports/final/../figures/M5_hypervolume.pdf}
\caption{Normalised hypervolume against evaluations for each method, mean and
min--max band over 5 seeds, Fidelity 1, run
\texttt{M5-\allowbreak{}ABL-\allowbreak{}20260920T211134Z}. File
\texttt{reports/\allowbreak{}figures/\allowbreak{}M5\_\allowbreak{}hypervolume.png}.
Per-seed curves:
\texttt{M5\_\allowbreak{}hypervolume\_\allowbreak{}per\_\allowbreak{}seed.png};
fronts: \texttt{M5\_\allowbreak{}fronts.png}; budget use:
\texttt{M5\_\allowbreak{}budget\_\allowbreak{}use.png}; surrogate
calibration:
\texttt{M5\_\allowbreak{}surrogate\_\allowbreak{}calibration.png}.}
\end{figure}

\subsection{What the agent did}

A hand-written audit of all 67 agent rounds
(\texttt{M5\_\allowbreak{}qualitative\_\allowbreak{}audit.md}) is the
only handle on the prior-knowledge question. Its findings:

The direction of every lever was stated in round 1, before any seed had
more than three feasible designs. 35 of 47 rounds say explicitly that
they rely on prior knowledge, 12 naming Allen--Eggers and 9
Sutton--Graves. First-round proposals had the direction right on 88 of
95 stated flux changes and 86 of 93 bondline changes. The agent starts
where the other optimisers have to get to, and its lead at 50
evaluations measures priors plus data. The study cannot separate the
two.

Some things came from the data and nowhere else. Three seeds derived in
round 1, from the validator's messages, a closed-form validity rule that
agrees with the evaluator on all 5400 paid designs of the run. Every
seed found the CFD-hull edge from extrapolation refusals by its second
or third call. Every seed stated within three to five rounds that the
front is a one-parameter curve in entry angle closed by the g-limit and
the mass fraction, which is §10.4's reading.

The agent is the most precise fence-parker in the study: 64\% of its
front designs have a mass-fraction margin below 10\textsuperscript{−4}.
Dropping every design within 0.1\% of any constraint reduces its final
lead over \texttt{bo\_\allowbreak{}parego} from 0.0036 to 0.0020. 195 of
its 1000 paid designs lie within 10\textsuperscript{−3} of an earlier
one in the unit cube; removing them lowers its mean hypervolume from
0.35557 to 0.35517. In 67 rounds it never asked whether a shield
weighing as much as the vehicle was a meaningful place to be (NR-32).

It also lost a correct rule. Seed 37 derived the validity rule in round
1; the prompt's digest of earlier rounds carries only the first 400
characters of one field; in round 2 the agent refitted a looser rule and
sent all 20 proposals into the invalid region. That is a property of the
harness.

Surrogate accuracy for \texttt{bo\_\allowbreak{}parego} was scored
prospectively on 900 logged predictions: 83\% of its picks lay outside
the training hull and were flagged as extrapolations when made; bondline
RMSE was 0.811 K inside the hull and 1.49 K outside, with intervals
somewhat overconfident (z-score spread 1.2 to 1.6).

The comparison holds for this model and this design space. The problem
has one real trade-off and monotone free improvements, the most
favourable case for prior knowledge, and it says nothing about a problem
where the textbook direction is wrong. One model, five seeds.

\section{Adaptive Fidelity}

Fidelity 1 with real OpenFOAM calls. Run
\texttt{M6-\allowbreak{}AF-\allowbreak{}20260920T211148Z}. Five arms
(adaptive promotion, random, greedy, up-front space-filling, and no CFD)
each ran NSGA-II for 200 cheap evaluations with a budget of 8
coarse-mesh CFD calls, over 5 seeds, all starting from the 69-point
surface. A usable CFD case joins that arm's training set and the surface
is refitted. Every arm was scored on a pooled-truth surface through all
178 CFD points, on the designs the arm's own final belief recommends,
with recommendations that truth calls infeasible dropped. The reference
hypervolume, 0.3436, comes from a dedicated 2 × 1000-evaluation search
on the truth surface.

The LLM-guided adaptive arm declared in the configuration was not run;
its live calls were not authorised for this run and it lies outside the
declared criterion. H2's ``AI-guided'' clause is therefore untested.

The verdict, by the pre-declared rule, is
\texttt{NOT\_\allowbreak{}SUPPORTED}. The adaptive arm reached the 95\%
target (0.3264) in 0 of 5 seeds; it needed 3. No arm reached it in any
seed. The verdict is not \texttt{NOT\_\allowbreak{}TESTABLE}, which is
reserved for the no-CFD arm reaching the target.

\begin{table}[htbp]
\caption{Arm scores (source:
\texttt{reports/\allowbreak{}milestones/\allowbreak{}M6\_\allowbreak{}adaptive\_\allowbreak{}fidelity.md}
§4; \texttt{M6\_\allowbreak{}addendum\_\allowbreak{}posthoc.md}).}
\small\centering
\begin{tabular}{@{}>{\raggedright\arraybackslash}p{\dimexpr(\linewidth-8\tabcolsep)*250/1000\relax}>{\raggedright\arraybackslash}p{\dimexpr(\linewidth-8\tabcolsep)*250/1000\relax}>{\raggedright\arraybackslash}p{\dimexpr(\linewidth-8\tabcolsep)*250/1000\relax}>{\raggedright\arraybackslash}p{\dimexpr(\linewidth-8\tabcolsep)*250/1000\relax}@{}}
\toprule
\tabhead{Arm} & \tabhead{Truth hypervolume, mean ± s.d.} & \tabhead{CFD calls (mean)} & \tabhead{Failed calls (mean)} \\
\midrule
Up-front & 0.2972 ± 0.0262 & 8.0 & 0.6 \\
Adaptive & 0.2835 ± 0.0389 & 3.6 & 0.4 \\
Greedy & 0.2823 ± 0.0395 & 7.2 & 0.8 \\
No CFD & 0.2754 ± 0.0375 & 0.0 & 0.0 \\
Random & 0.2652 ± 0.0254 & 8.0 & 0.8 \\
\bottomrule
\end{tabular}
\end{table}

No paired comparison of the adaptive arm with any other is
distinguishable from zero at n = 5 (p from 0.11 to 0.73). With five
seeds that is weak evidence of equivalence.

Three facts explain the result. First, extra CFD did not change what the
designs are worth: the no-CFD arm's recommended sets score 0.2753 on the
starting surface and 0.2754 on pooled truth, so 109 additional CFD
points moved it by 0.0001. No arm made a false feasibility claim, and
the largest optimism gap is 0.0004. The starting surface was already
accurate where the objectives are decided, consistent with C\_D at peak
heating varying by about 0.1\% across the front.

Second, the target was out of reach of the study's own search budget, a
sizing flaw (NR-35). All 25 arm-seeds pooled, 5000 evaluations, reach
94.1\% of the reference, below the 95\% target. A saving in CFD calls
could not have shown up in this criterion. The configuration's sizing
had checked statistical power and machine time, not reachability.

Third, the searches never reached the front's region. Reference-front
shapes have R\_n/D from 0.999 to 1.133; no arm found a feasible design
above 0.984 in 200 evaluations, and none of the promoting arms' 94 calls
was at R\_n/D ≥ 1.0. Seed explains 68.7\% of the variance in final
hypervolume and arm 10.6\%.

The adaptive arm matched the greedy arm's score on half the calls. That
is true, and it is not the criterion, so it is not claimed as support.
11 of 125 CFD cases failed and stay on record, charged.

In this design space at this budget, extra CFD bought nothing
measurable, so adaptive allocation had nothing to save. H2 is not
supported. It is not contradicted in general either: a space where the
cheap model is materially wrong near the front is where H2 would have
something to show, and this study as sized could not have detected it. A
reachable criterion would be a new pre-declared study. I have left this
one as written and not re-scored it.

\begin{figure}[htbp]
\centering
\includegraphics[width=\linewidth,height=0.42\textheight,keepaspectratio]{/Users/akj/Projects/aether/reports/final/../figures/M6_cfd_spend_map.pdf}
\caption{Where each arm spent its CFD calls in bluntness ratio and cone
half-angle (degrees), against the reference-front shapes, run
\texttt{M6-\allowbreak{}AF-\allowbreak{}20260920T211148Z} (shoulder
ratio and Mach not shown). File
\texttt{reports/\allowbreak{}figures/\allowbreak{}M6\_\allowbreak{}cfd\_\allowbreak{}spend\_\allowbreak{}map.png}.
Truth hypervolume against CFD calls:
\texttt{M6\_\allowbreak{}hv\_\allowbreak{}vs\_\allowbreak{}cfd\_\allowbreak{}calls.png}
(curves rise with calls partly because search progresses while calls are
spent); hold-out error of the truth surface:
\texttt{M6\_\allowbreak{}surface\_\allowbreak{}error\_\allowbreak{}vs\_\allowbreak{}truth.png}.}
\end{figure}

\section{Uncertainty}

Fidelity 1. Propagation and attribution run
\texttt{M7-\allowbreak{}UQ-\allowbreak{}20260920T233048Z}; robust
optimisation run
\texttt{M7-\allowbreak{}ROBUST-\allowbreak{}20260920T211216Z};
like-for-like propagation of the robust knee
\texttt{M7-\allowbreak{}LFL-\allowbreak{}20260921T011854Z}.

\subsection{Inventory and sampling}

Twelve active inputs: three aleatory (atmospheric density, vehicle mass
±2\% 1σ, delivered entry angle) and nine epistemic. The epistemic set
includes a discrete 50/50 switch between the two effective-nose-radius
primaries, the Sutton--Graves constant ±4\%, a 0.5 to 1.5 multiplier on
heating above 86 km, TPS conductivity ±15\%, specific heat ±10\%, and
four drag terms (GP, discretisation, the perfect-gas ±5\% band, base
drag). Seven of the twelve are engineering judgment, including both TPS
property bands and the high-altitude multiplier. Density dispersions
below 86 km are NASA-published standard deviations {[}28{]}; those above
are this project's conversion of observed ranges.

\begin{figure}[htbp]
\centering
\includegraphics[width=\linewidth,height=0.42\textheight,keepaspectratio]{/Users/akj/Projects/aether/reports/final/../figures/infographics/05_uncertainty_sources.pdf}
\caption{Where the uncertainty comes from: 84 to 99\% of the output variance is
epistemic, and the two dominant bondline terms (TPS conductivity,
heating above 86 km) are engineering judgment. Explanatory figure; tier
and aleatory/epistemic labels are the inventory's own fields, read from
\texttt{results/\allowbreak{}M7/\allowbreak{}M7-\allowbreak{}UQ-\allowbreak{}20260920T233048Z/\allowbreak{}summary.json}
(attribution, propagation decomposition, uncertainty model). Every
spread shown is a lower bound (§13.1). File
\texttt{reports/\allowbreak{}figures/\allowbreak{}infographics/\allowbreak{}05\_\allowbreak{}uncertainty\_\allowbreak{}sources.png}.}
\end{figure}

Sampling is nested: 24 epistemic branches × 125 aleatory draws with
common random numbers, 3000 full evaluations per design. Aleatory and
epistemic uncertainty are kept apart by the sample shape, so the primary
output is a band of distributions across branches.

This is an uncertainty study of a model. It cannot cover model-form
error that is not in the inventory, above all G2{\symfont ′}. Every
spread below is a lower bound.

\subsection{Propagation}

\begin{table}[htbp]
\caption{Peak bondline temperature under uncertainty, K (source:
\texttt{reports/\allowbreak{}milestones/\allowbreak{}M7\_\allowbreak{}uncertainty\_\allowbreak{}robust.md}
§3.1.2).}
\small\centering
\begin{tabular}{@{}>{\raggedright\arraybackslash}p{\dimexpr(\linewidth-12\tabcolsep)*232/1000\relax}>{\raggedright\arraybackslash}p{\dimexpr(\linewidth-12\tabcolsep)*136/1000\relax}>{\raggedright\arraybackslash}p{\dimexpr(\linewidth-12\tabcolsep)*109/1000\relax}>{\raggedright\arraybackslash}p{\dimexpr(\linewidth-12\tabcolsep)*109/1000\relax}>{\raggedright\arraybackslash}p{\dimexpr(\linewidth-12\tabcolsep)*109/1000\relax}>{\raggedright\arraybackslash}p{\dimexpr(\linewidth-12\tabcolsep)*301/1000\relax}@{}}
\toprule
\tabhead{Design} & \tabhead{Nominal} & \tabhead{Mean} & \tabhead{s.d.} & \tabhead{p95} & \tabhead{p95 band over branches} \\
\midrule
Baseline & 504.3 & 503.4 & 13.75 & 523.1 & {[}478.7, 526.2{]} \\
Peak-flux-only & 418.4 & 415.1 & 7.632 & 426.2 & {[}404.2, 430.3{]} \\
Joint knee & 403.5 & 400.8 & 6.464 & 410.4 & {[}391.2, 412.7{]} \\
Bondline-only & 391.7 & 389.5 & 5.638 & 398.0 & {[}380.8, 399.1{]} \\
\bottomrule
\end{tabular}
\end{table}

The epistemic share of variance is 84.4 to 93.9\% for peak flux and 97.2
to 98.7\% for bondline temperature. Ignorance of this kind is removed by
evidence, not by design margin.

The pre-declared convergence check, a bootstrap over rows, is met by all
eight statistics. It understates the sampling error: rows within a
branch are not independent, and resampling whole branches gives
half-widths 9 to 18 times larger. On that stricter reading the baseline
mean marginally misses the 1\% tolerance (1.06\%). Absolute means are
known to a few kelvin. More branches would tighten them; more draws
would not (NR-34).

\begin{figure}[htbp]
\centering
\includegraphics[width=\linewidth,height=0.42\textheight,keepaspectratio]{/Users/akj/Projects/aether/reports/final/../figures/M7_pbox.pdf}
\caption{Probability boxes for peak bondline temperature (K): one empirical
distribution per epistemic branch for each design, 24 branches × 125
draws, Fidelity 1, run
\texttt{M7-\allowbreak{}UQ-\allowbreak{}20260920T233048Z}. File
\texttt{reports/\allowbreak{}figures/\allowbreak{}M7\_\allowbreak{}pbox.png}.
Input inventory:
\texttt{M7\_\allowbreak{}input\_\allowbreak{}inventory.png}; pooled
output distributions:
\texttt{M7\_\allowbreak{}output\_\allowbreak{}distributions.png};
convergence: \texttt{M7\_\allowbreak{}convergence.png}.}
\end{figure}

\subsection{Is the knee's advantage larger than the uncertainty on it?}

Each design's own s.d. is 6.5 to 7.6 K, so read as independent clouds
the 14.9 K difference would look marginal. That reading is wrong,
because all designs were propagated on one shared draw set: draw
\emph{i} is the same atmosphere, mass, conductivity and nose-radius
model for every design.

\begin{table}[htbp]
\caption{Paired same-draw differences in peak bondline temperature, K (source:
\texttt{M7\_\allowbreak{}addendum\_\allowbreak{}posthoc.md} §1;
\texttt{paired\_\allowbreak{}difference.json}).}
\small\centering
\begin{tabular}{@{}>{\raggedright\arraybackslash}p{\dimexpr(\linewidth-12\tabcolsep)*264/1000\relax}>{\raggedright\arraybackslash}p{\dimexpr(\linewidth-12\tabcolsep)*90/1000\relax}>{\raggedright\arraybackslash}p{\dimexpr(\linewidth-12\tabcolsep)*66/1000\relax}>{\raggedright\arraybackslash}p{\dimexpr(\linewidth-12\tabcolsep)*165/1000\relax}>{\raggedright\arraybackslash}p{\dimexpr(\linewidth-12\tabcolsep)*140/1000\relax}>{\raggedright\arraybackslash}p{\dimexpr(\linewidth-12\tabcolsep)*272/1000\relax}@{}}
\toprule
\tabhead{Difference} & \tabhead{Mean} & \tabhead{s.d.} & \tabhead{Range} & \tabhead{Draws with Δ \textless{} 0} & \tabhead{95\% CI on mean (branch bootstrap)} \\
\midrule
Knee − peak-flux-only & −14.31 & 1.45 & −17.85 to −10.80 & 3000 / 3000 & {[}−14.88, −13.76{]} \\
Bondline-only − peak-flux-only & −25.64 & 2.44 & −31.55 to −19.79 & 3000 / 3000 & {[}−26.59, −24.73{]} \\
\bottomrule
\end{tabular}
\end{table}

The branch-mean difference for the knee lies in {[}−17.05, −11.41{]} K
and is negative in all 24 branches. Which model is right changes the
size of the advantage by about ±3 K, not its sign. The knee's peak flux
is higher in all 3000 draws, by 20.5 kW/m\textsuperscript{2} on average.
Zero of 3000 is an upper bound of 0.13\% at 95\% confidence, within the
declared uncertainty model. A model error common to all designs moves
them together, and a paired test cannot see it.

\subsection{What drives the spread}

\begin{table}[htbp]
\caption{Largest total-order Sobol' indices (source: M7 report §4).}
\small\centering
\begin{tabular}{@{}>{\raggedright\arraybackslash}p{\dimexpr(\linewidth-6\tabcolsep)*142/1000\relax}>{\raggedright\arraybackslash}p{\dimexpr(\linewidth-6\tabcolsep)*428/1000\relax}>{\raggedright\arraybackslash}p{\dimexpr(\linewidth-6\tabcolsep)*428/1000\relax}@{}}
\toprule
\tabhead{Output} & \tabhead{Baseline (K = 1, hemisphere)} & \tabhead{Joint knee (K = 0.417)} \\
\midrule
Peak heat flux & Sutton--Graves constant 0.609; perfect-gas C\_D form 0.258; mass 0.122;
nose-radius model 0.000 & nose-radius model 0.706 {[}0.599, 0.819{]}; Sutton--Graves 0.181;
perfect-gas form 0.065 \\
Peak bondline T & TPS conductivity 0.883; heating above 86 km 0.086 & TPS conductivity 0.638; heating above 86 km 0.262; nose-radius model
0.052 \\
Max g & entry angle 0.340; density 0.285; base drag 0.155; perfect-gas form
0.154 & entry angle 0.487; density 0.338; base drag 0.121 \\
\bottomrule
\end{tabular}
\end{table}

At the front, the disagreement between the two NASA primaries carries
71\% of the peak-flux variance. It is exactly zero on the hemispherical
baseline, where the two agree identically. It is a term that evidence
could resolve. The bondline is governed by the insulator's conductivity
and by heating above 86 km, both engineering judgment. The second
objective is dominated by the two numbers the project has least evidence
for. The GP surrogate term (≤ 0.020) and the mesh term (≤ 0.003) are
negligible for every output.

\begin{figure}[htbp]
\centering
\includegraphics[width=\linewidth,height=0.42\textheight,keepaspectratio]{/Users/akj/Projects/aether/reports/final/../figures/M7_attribution.pdf}
\caption{Total-order Sobol' indices over the uncertain inputs for peak heat flux,
peak bondline temperature and maximum deceleration, designs
\texttt{baseline} and \texttt{joint\_\allowbreak{}knee}, run
\texttt{M7-\allowbreak{}UQ-\allowbreak{}20260920T233048Z} (regenerated
2026-09-21; an earlier version of this figure was mislabelled, NR-34).
File
\texttt{reports/\allowbreak{}figures/\allowbreak{}M7\_\allowbreak{}attribution.png}.}
\end{figure}

\subsection{Fragility of the nominal optima, and the chance-constrained
front}

None of the three nominal optima is chance-feasible at the declared 5\%.
P(any constraint violated) is 34.40\% {[}32.72, 36.12{]} for the
peak-flux-only design, 29.60\% {[}27.99, 31.26{]} for the knee and
46.60\% {[}44.82, 48.39{]} for the bondline-only design. None of it is
thermal: the bondline constraint is violated in 0 of 3000 draws for
each. The crossed constraint is the mass-fraction fence, pushed over by
the ±2\% mass dispersion on a 350 kg placeholder, plus the 12 g limit
for the bondline-only design (36.00\%). The designs sit 0.84\%, 1.09\%
and 0.19\% from their binding constraint at nominal, which is where an
optimiser leaves them. The fragility is a fact about this model's
constraint set, two placeholders and a fence.

Robust optimisation minimised the 95th percentiles of both objectives
under chance constraints at 5\%, with 32 inner draws per candidate under
common random numbers, 3 seeds × 30,000 inner evaluations. Eight front
designs were re-scored with 1000 fresh independent draws. The
pre-declared shortcut verification passed: relative bias −1.61\% and
+0.32\%, Spearman 1.000 on both objectives, largest probability error
2.60\%.

Two things sit beside that pass. All 46 robust-front designs show a
mass-fraction violation probability of exactly 1/32 on the inner sample:
the optimiser used its allowance to the last draw on 32 draws it could
learn. And on fresh draws one of the eight checked designs, at the
low-bondline end, reads P(any) = 5.7\% against the 5\% limit. A flipped
verdict is not one of the three declared tolerances, so the check
passes, and the flip is reported as its own line (NR-33).

The trade-off survives among chance-feasible designs: the robust front
spans 30.6 K of 95th-percentile bondline temperature, and its
flux-minimising end is its hottest-bondline end. At the knee, robustness
costs +7.8 kW/m\textsuperscript{2} (+3.1\%) of nominal peak flux,
through a diameter 1.0 to 1.5\% smaller, and brings P(any violation)
from 29.60\% to 3.20\% {[}2.63, 3.89{]} on the same 3000 draws. Paired
with the nominal knee, its bondline is lower by 0.42 K (s.d. 0.12 K) in
all 3000 draws.

\begin{figure}[htbp]
\centering
\includegraphics[width=\linewidth,height=0.46\textheight,keepaspectratio]{/Users/akj/Projects/aether/reports/final/../figures/M7_robust_vs_nominal.pdf}
\caption{Nominal Pareto front (78 designs, run
\texttt{M4-\allowbreak{}OPT-\allowbreak{}20260920T205252Z}) and
chance-constrained robust front (95th percentiles, run
\texttt{M7-\allowbreak{}ROBUST-\allowbreak{}20260920T211216Z}), peak
heat flux (W/m\textsuperscript{2}) against peak bondline temperature
(K); arrows join each nominal design to its own 95th percentile. File
\texttt{reports/\allowbreak{}figures/\allowbreak{}M7\_\allowbreak{}robust\_\allowbreak{}vs\_\allowbreak{}nominal.png}.
Shortcut verification:
\texttt{M7\_\allowbreak{}shortcut\_\allowbreak{}verification.png}.}
\end{figure}

\subsection{Final comparison}

\begin{table}[htbp]
\caption{Final comparison. All four designs were re-evaluated under one evaluator
source hash; Fidelity 1
(\texttt{cfd\_\allowbreak{}surface\_\allowbreak{}v2}); uncertainty rows
from 3000 nested draws, seed 20260921 (source: M7 report §6;
\texttt{comparison.csv}).}
\small\centering
\begin{tabular}{@{}>{\raggedright\arraybackslash}p{\dimexpr(\linewidth-10\tabcolsep)*268/1000\relax}>{\raggedright\arraybackslash}p{\dimexpr(\linewidth-10\tabcolsep)*171/1000\relax}>{\raggedright\arraybackslash}p{\dimexpr(\linewidth-10\tabcolsep)*179/1000\relax}>{\raggedright\arraybackslash}p{\dimexpr(\linewidth-10\tabcolsep)*186/1000\relax}>{\raggedright\arraybackslash}p{\dimexpr(\linewidth-10\tabcolsep)*194/1000\relax}@{}}
\toprule
\tabhead{Metric} & \tabhead{Baseline} & \tabhead{Peak-heat-only optimised} & \tabhead{Joint O1 optimised (knee)} & \tabhead{Robust O1 optimised (knee)} \\
\midrule
Peak heat flux {[}W/m\textsuperscript{2}{]} & 1.868×10\textsuperscript{6} & 2.295×10\textsuperscript{5} & 2.510×10\textsuperscript{5} & 2.588×10\textsuperscript{5} \\
Integrated external heat {[}J/m\textsuperscript{2}{]} & 1.306×10\textsuperscript{8} & 1.975×10\textsuperscript{7} & 1.781×10\textsuperscript{7} & 1.801×10\textsuperscript{7} \\
Peak surface temperature {[}K{]} & 2476 & 1428 & 1459 & 1471 \\
Peak bondline temperature {[}K{]} & 504.3 & 418.4 & 403.5 & 403.0 \\
Thermal penetration depth {[}m{]} & 0.02475 & 0.01281 & 0.01256 & 0.01256 \\
TPI {[}K·m·s{]}, for completeness only & 2875 & 1011 & 831.6 & 828 \\
Max g & 12.25 & 9.216 & 10.35 & 10.48 \\
Max dynamic pressure {[}Pa{]} & 4.184×10\textsuperscript{4} & 2618 & 2940 & 3065 \\
Entry duration {[}s{]} & 240 & 390.7 & 340.1 & 333.9 \\
Feasible at nominal & no & yes & yes & yes \\
p95 peak heat flux {[}W/m\textsuperscript{2}{]} & 1.952×10\textsuperscript{6} & 2.352×10\textsuperscript{5} & 2.571×10\textsuperscript{5} & 2.651×10\textsuperscript{5} † \\
p95 peak bondline T {[}K{]} & 523.1 & 426.2 & 410.4 & 410.0 † \\
P(any constraint violated) & 100.00\% {[}99.87, 100.00{]} & 34.40\% {[}32.72, 36.12{]} & 29.60\% {[}27.99, 31.26{]} & 3.20\% {[}2.63, 3.89{]} † \\
\bottomrule
\end{tabular}
\end{table}

† Like-for-like cell from run
\texttt{M7-\allowbreak{}LFL-\allowbreak{}20260921T011854Z}: the robust
knee propagated through the same 3000 nested draws as the other columns,
under a later source hash that first reproduced 120 of the study's own
evaluations to 3.5×10\textsuperscript{−14} relative. As the study
originally generated it, this cell was n = 500 mixed draws on a
different seed: p95 flux 2.699×10\textsuperscript{5}
W/m\textsuperscript{2}, p95 bondline 414.1 K, P(violation) 2.80\%
{[}1.68, 4.64{]} (14 of 500). That cell is kept in the milestone report
and superseded here only because it could not be paired with the other
columns.

``Feasible'' is measured against a 12 g limit that is an open decision,
a 450 K substrate-specific allowable, and a mass-fraction fence. Under
option A of the deceleration decision it means survivable in a
contingency, not acceptable for a nominal crewed entry. The TPI row is
printed so a reader can check the redundancy claim of §17.2 on these
designs; it ranks nothing.

\section{Physical Experiment}

\textbf{The experiment has not been run.} No coupon has been printed, no
heater switched on, and no measurement exists anywhere in the
repository. This section describes the designed experiment and the
blind-prediction protocol, and leaves a marked slot for results.

\subsection{Purpose and design}

The headline result rests on one piece of physics: a slab of solid acts
as a diffusive low-pass filter on a heat pulse, so interior temperature
history depends on the pulse's shape as well as its energy. That has
been checked only against an analytical solution, which verifies
arithmetic. The experiment heats a 60 × 60 × 10 mm solid-printed polymer
coupon, instrumented at four or more depths plus ambient, with two
heating histories carrying the same integrated energy, one high-peak and
short, one lower-peak and long, and asks whether the conduction model,
calibrated on different runs and then frozen, predicts both in-depth
histories
(\texttt{experiments/\allowbreak{}thermal\_\allowbreak{}coupon/\allowbreak{}README.md}).

It is a thermal-transient validation experiment. It is not a re-entry
simulator: a hobby heater on a plastic plate near 350 K shares nothing
with a 2 MW/m\textsuperscript{2} stagnation point. It is also not a test
of the burn-versus-bake conclusion. In §7 the longer pulse also carries
a larger integrated load; here energy is held fixed on purpose, which
isolates duration and, according to the package's own design
calculation, reverses the ordering. The predicted difference between the
two cases at the deepest sensor is a few kelvin (about 3.5 K for the
placeholder properties), so sensor resolution and repeatability matter
more than range. The fitted properties will be effective
through-thickness values for one printed part, not datasheet values.

\subsection{Blind-prediction protocol}

\begin{enumerate}
\def\labelenumi{\arabic{enumi}.}
\tightlist
\item
  Two-point sensor calibration and a calorimetric calibration of heater
  flux.
\item
  A dedicated steady-state run to measure heater-to-coupon contact
  resistance, which is then held fixed.
\item
  Two calibration runs: a cool-down, which pins the loss coefficients,
  and a constant-flux step held for at least two to three diffusion
  times, about an hour for a 10 mm coupon, which separates conductivity
  from specific heat.
\item
  Declare, in a committed notebook entry, the RMSE, peak error and
  time-to-peak error that will count as agreement, before any validation
  run. The comparison script assigns no verdict.
\item
  Freeze the parameters. Generate predictions for both validation cases.
  The prediction script refuses to run if the measurement file already
  exists. Commit and push predictions and manifests.
\item
  Run each validation case at least twice, logging through the
  post-heating soak, then compare on RMSE, peak error, time-to-peak
  error, residual structure and uncertainty coverage.
\end{enumerate}

Steps 2 and 3, and the choice of peak estimator, came from exercising
the whole chain on synthetic data before anything was built (§17.5).
With those fixes the synthetic chain recovers conductivity, specific
heat and both loss coefficients to better than 0.5\%, fitting on a
coarser mesh than the one that generated the data, and predicts an
unseen pulse to better than 0.1 K against noise-free truth. That is a
statement about the software. No real material was involved.

Still open and the author's to decide: which polymer and therefore the
temperature ceiling, which heater, thermocouples or thermistors, how
flux is known, and the tolerance of step 4.

\subsection{Results}

\begin{quote}
\textbf{{[}PENDING --- STUDENT{]} Results slot.} To be filled only after
the protocol has been run: declared tolerance and its commit hash;
frozen parameters; archived prediction manifest hashes; for each
validation case and sensor, RMSE, peak error (smoothed-plateau
estimator), time-to-peak error, residual plots and uncertainty coverage;
whether the two cases differ in the predicted direction and by the
predicted magnitude; anything that went wrong.
\end{quote}

\section{Discussion}

The mechanism is verified and the magnitudes are model-dependent. That
separation runs through every result above.

The mechanism is a comparison of two timescales inside a conduction
solver that matches an analytical solution to 0.002\%. Within the model
it holds across every perturbation tried: across the one-parameter
sweep, across a change of drag model that moved the baseline's peak flux
by 16.52\%, across a correction to the heating model worth 20\% on the
optimised designs, and across 3000 draws and 24 epistemic branches of
the declared uncertainties. The ordering of the three front designs is
not an artefact of the nominal point.

The magnitudes are less secure. The 14.9 K at the knee is a temperature
difference inside a placeholder stack, at a geometry selected by a mass
fence, a CFD hull edge, a box bound and a frozen shoulder. The
Fidelity-1 drag model was built so that shape could matter, and the
front it produced has the same shape everywhere. H1 is supported by the
same physics as H0, entry steepness, and says nothing yet about capsule
shape. A study in which shape could trade the objectives needs a mass
model in which vehicle mass grows with size, a heating model that can
see the shoulder, and CFD anchors beyond R\_n/D = 1.2.

What would change the conclusions, in order of likely size. A catalycity
or hot-wall treatment would move every heat flux, possibly by a factor
near two, but in common mode; it would change levels and allowable
margins more than orderings. An ablating TPS would carry energy away
disproportionately in the long shallow entry and shrink the effect. A
stack with much more insulation would remove the effect, as it did at 40
mm. A sourced mass budget would move the front to smaller diameters and
higher heating throughout. Adopting the documented deceleration curve
would empty the Fidelity-0 feasible region, and its effect on the
Fidelity-1 front has not been measured because the evaluator stores peak
g only. And a finding that the conduction model mispredicts a real
coupon would undercut the one \texttt{PASS} the result rests on.

The uncertainty attribution reads as a list of evidence to obtain. For
peak flux at the front, it is the effective nose radius at K ≈ 0.42: the
single highest-value follow-up has a title and a DOI (10.2514/1.T7458)
and is paywalled. For the bondline, it is a measured conductivity for an
actual insulator and a defensible treatment of heating above 86 km,
where a continuum correlation is being applied to transitional flow and
an interpolated atmosphere.

On the two method questions. The language-model agent reached the
front's structure sooner than the numeric optimisers and parked closer
to two placeholder limits; nothing shows it found a design or lever that
NSGA-II with five times the budget did not. The adaptive-fidelity study
found that extra CFD bought nothing here, because §9's surface was
already accurate where the objectives are decided. Read with §13.4,
where the surrogate and mesh terms carry at most 0.02 of any output's
variance, the second finding says that the drag model is not where this
project's uncertainty lies.

\section{Limitations}

Assembled from \texttt{ASSUMPTIONS.md} and the limitation sections of
the milestone reports.

\subsection{Fidelity of reported numbers}

§7 is Fidelity 0 with the legacy nose-radius setting (harmless at K = 1)
and a constant C\_D that §9 shows to be wrong by +16.52\% in peak flux
for the baseline. §10 to §13 are Fidelity 1, which means a better drag
model, not a validated one.

\subsection{Gate statuses}

\texttt{LIMITED}: G1A{\symfont ′} outside five altitudes; G2{\symfont ′}
entirely; G-GEO2 physics; the drag surface (no capsule validation data,
coarse mesh, sphere-derived discretisation band, perfect gas, assumed
base drag, per-case iterative scatter up to about 0.2\%); every study
row M4 to M7 (no external reference exists for any of them). G4 is
\texttt{PASS} under a restart rule written after the first results were
seen. G5 is a software gate. M8 is \texttt{IN\_\allowbreak{}PROGRESS}.

\subsection{Heating, aerodynamics and TPS}

Heating is stagnation point only, convective only, cold wall, catalycity
unquantified. Heating above 86 km, 5 to 21\% of integrated load for
feasible and front designs, lies outside the continuum regime
Sutton--Graves describes. The effective-nose-radius model rests on Mach
8 tunnel data transferred by a pressure-invariance argument, on this
project's own interpolation between nine points, and on two primaries
that disagree by about 20\% where the front sits.

The aerodynamics are inviscid, perfect gas, zero angle of attack,
forebody only. The ±5\% model-form band is declared and has not been
validated. No lift, trim or stability exists in the model, so the
stability constraint named in the objective is not checked at all.

The TPS model has placeholder properties, constant with temperature, no
ablation, adiabatic back face, no contact resistance, one-dimensional.

\subsection{Constraints}

The deceleration limit is an open decision for the author. Flat 12 g
appears in no NASA document found. Option A keeps 12 g and relabels it
an emergency envelope; option B adopts the NASA-STD-3001 deconditioned
duration curve, which needs a pulse-duration metric the evaluator does
not compute and which, measured on the Fidelity-0 grid, leaves 4 of 140
designs feasible at 10 g and none at 8 g. The mass-fraction limit of 1.0
is a fence; no sourced mass budget exists, and the knee's shield is 346
kg of a 350 kg vehicle. Variable ranges are engineering choices.

\subsection{Statistics and process}

Five seeds for the optimiser comparison and the adaptive-fidelity study;
seven for §10 with no significance test. Twenty-four epistemic branches.
Seven of twelve uncertainty inputs are engineering judgment. Every
uncertainty figure is a lower bound.

Most of the code was AI-generated and most of it has not yet been read
line by line by the author (§19). No external technical review has been
requested or received. The physical experiment has not been run. Several
milestone reports carry a dirty-working-tree flag in their headers.

\section{Negative Results}

Thirty-five entries are kept in
\texttt{docs/\allowbreak{}negative\_\allowbreak{}results.md}, including
ones whose evidence is a scratch run that no longer exists, labelled as
such. This section reports those that changed what the project claims.
It sits before the conclusion because several numbers above cannot be
read correctly without it.

\begin{figure}[htbp]
\centering
\includegraphics[width=\linewidth,height=0.42\textheight,keepaspectratio]{/Users/akj/Projects/aether/reports/final/../figures/infographics/06_negative_results_timeline.pdf}
\caption{The negative-results timeline: the thirty-five entries of
\texttt{docs/\allowbreak{}negative\_\allowbreak{}results.md} in four
kinds, with five called out. Explanatory figure; titles are the file's
verbatim headings, dates are from \texttt{git log -\allowbreak{}S} on
that file, and the grouping by kind is the rendering script's own
classification, labelled as such on the figure. File
\texttt{reports/\allowbreak{}figures/\allowbreak{}infographics/\allowbreak{}06\_\allowbreak{}negative\_\allowbreak{}results\_\allowbreak{}timeline.png}.}
\end{figure}

\subsection{The model had holes, and the optimiser found them}

\subsubsection{A heat shield heavier than the vehicle
(NR-13)}\label{a-heat-shield-heavier-than-the-vehicle-nr-13}

Drag area grows with D\textsuperscript{2} while vehicle mass was a fixed
configuration value, so both objectives improve with diameter and
nothing charged for it. I had anticipated this by arithmetic, and the
first Latin Hypercube confirmed it. Under the two original constraints
the non-dominated set was four designs at 3.9 to 4.5 m whose forebody
TPS alone weighed 1.29 to 2.81 times the entire vehicle, and 482 of 1328
evaluable samples violated mass closure. A constraint at 1.0 excludes
only the impossible, and the front then sits on it (§10.3). It is the
earlier finding that both Fidelity-0 optima sat on the diameter bound,
seen from the other side.

\subsubsection{An unbounded nose-flattening lever
(NR-15)}\label{an-unbounded-nose-flattening-lever-nr-15}

With R\_eff taken as the cap radius, flattening the nose lowered both
objectives and cost nothing. All 64 designs on the Fidelity-0 front
ended at a bluntness ratio of 1.18 to 1.20 against a cap of 1.2. The cap
had been placed in advance, and recorded at the time as the wrong kind
of fix: a constraint at the edge of validity is a fence, not a model,
and ``the optimiser chose 1.2'' meant ``the fence is at 1.2''. What the
fence was hiding, measured once the corrected model existed:

\begin{center}
\small\centering
\begin{tabular}{@{}>{\raggedright\arraybackslash}p{\dimexpr(\linewidth-6\tabcolsep)*333/1000\relax}>{\raggedright\arraybackslash}p{\dimexpr(\linewidth-6\tabcolsep)*333/1000\relax}>{\raggedright\arraybackslash}p{\dimexpr(\linewidth-6\tabcolsep)*333/1000\relax}@{}}
\toprule
\tabhead{Bluntness change} & \tabhead{Δ peak flux, legacy R\_eff = R\_n} & \tabhead{Δ peak flux, corrected} \\
\midrule
1.2 {\symfont →} 1.26 (where geometry refuses) & −2.41\% & −0.92\% \\
1.2 {\symfont →} 1.35 (box edge) & −5.72\% & −2.15\% \\
1.2 {\symfont →} infinitely flat & −98.90\% & −21.31\% \\
\bottomrule
\end{tabular}
\end{center}

An infinitely flat nose removing 98.9\% of stagnation heating is wrong
by construction, because 1/√R\_n goes to zero while a flat face has a
finite velocity gradient. Re-evaluating the 64 designs under both
models: R\_eff is 0.685 to 0.693 of the cap radius, peak flux rises 20.2
to 20.8\%, bondline temperature by 6.7 to 9.1 K, and all stay feasible.
All 64 landed essentially on one of Ellison's nine tabulated bodies, so
the correction interpolates almost nothing, which was luck. The stated
uncertainty got worse, because the primaries disagree by about 20\% at
exactly that geometry.

\subsubsection{The fix moved the exploit (NR-15 follow-up, NR-29,
NR-30)}\label{the-fix-moved-the-exploit-nr-15-follow-up-nr-29-nr-30}

Under the corrected model a sharper shoulder lowers stagnation heating,
and the model cannot see shoulder heating. I wrote that down before the
next run. The optimiser then never touched the shoulder, because the
screening rule froze it, and the probes of §10.3 show the frozen lever
is worth about 7\% at the front. The bluntness cap, removed from the
configuration, came back as the edge of where CFD had been run. Three
fences in sequence, each standing where the previous one stood, each
labelled.

\subsubsection{``Negligible above 86 km'' was false for the bondline
(NR-14)}\label{negligible-above-86-km-was-false-for-the-bondline-nr-14}

Over the first Latin Hypercube the share of integrated heat load
accruing above 86 km had a median of 9.5\% and a maximum of 27\%, and
among feasible designs was never below 5.0\%. It rises with diameter
(correlation 0.83), so the designs both objectives favour spend longest
where the atmosphere is an interpolated table and the heating
correlation is out of regime. Nothing was tuned; the share became a
per-candidate diagnostic.

\subsection{A metric proposed, tested against a threshold declared
first, and discarded (NR-10)}

The specification asked for a Thermal Penetration Index. It was
implemented, verified on closed-form cases, and evaluated over the
140-design grid against a redundancy criterion written into a
configuration file before the run (run
\texttt{TPI-\allowbreak{}20260920T124732Z}). Its ranking is 99.94\%
reconstructible by rank regression on peak bondline temperature and
integrated load; Spearman ρ against bondline temperature alone is
+0.9992; 91 of 9730 ordered pairs are ranked differently; a
TPI-minimising optimiser selects the same design as a
bondline-minimising one. All 16 combinations of reference temperature
and weighting fail together. Verdict: discard.

The pre-registered prediction was half wrong. The theory note predicted
the outcome and gave a mechanism, that the integrand would be dominated
by the outer millimetre or two. The outermost 20\% of depth carries only
27.6 to 40.6\% of the integral. The real mechanism is shape invariance:
normalised exceedance profiles of all 140 designs have a minimum
pairwise cosine similarity of 0.9366, so any weighted depth integral
recovers one scale factor. The distinction changes what future work
should try: the stack would have to change, not the weighting. The
prediction text is left as written with the addendum beside it.

\begin{figure}[htbp]
\centering
\includegraphics[width=\linewidth,height=0.42\textheight,keepaspectratio]{/Users/akj/Projects/aether/reports/final/../figures/TPI_weighting_and_profile.pdf}
\caption{Depth weightings and normalised time-integrated temperature-exceedance
profiles for the 140 grid designs, reference configuration T\_ref = 400
K, insulator depth 15 mm, run
\texttt{TPI-\allowbreak{}20260920T124732Z}. File
\texttt{reports/\allowbreak{}figures/\allowbreak{}TPI\_\allowbreak{}weighting\_\allowbreak{}and\_\allowbreak{}profile.png}.
TPI against bondline temperature:
\texttt{TPI\_\allowbreak{}vs\_\allowbreak{}bondline.png}; sensitivity to
the 16 configurations: \texttt{TPI\_\allowbreak{}sensitivity.png}; rank
residuals: \texttt{TPI\_\allowbreak{}rank\_\allowbreak{}residual.png}.}
\end{figure}

\subsection{Tools that produced a plausible number instead of an error}

\subsubsection{The Pareto front was wrong and a figure caught it
(NR-04)}\label{the-pareto-front-was-wrong-and-a-figure-caught-it-nr-04}

The first dominance test had inverted control flow and returned all 21
feasible designs as non-dominated. The plotted front zigzagged. A
two-objective front must be monotone, so the plot falsified the code. No
test had asserted monotonicity.

\subsubsection{A variance estimator gave impossible intervals
(NR-16)}\label{a-variance-estimator-gave-impossible-intervals-nr-16}

First-order Sobol' intervals on the bondline came out as wide as
{[}−0.41, 1.15{]}. The estimator's variance was dominated by the
output's 400 K mean. Centring the output, on the same 6144 evaluations,
gave {[}0.30, 0.45{]}. Total-order indices were unaffected, so no
screening decision changed.

\subsubsection{A Bayesian optimiser quietly stopped optimising
(NR-17)}\label{a-bayesian-optimiser-quietly-stopped-optimising-nr-17}

A GP classifier returned NaN, the acquisition became all-NaN, and the
loop fell through to space-filling samples. Hypervolume did not show it,
because on this problem the front is found early. It showed in where the
budget went. A second fault sent 60\% of the budget to capsules that
cannot exist. These observations came from a scratch run no longer on
disk and are recorded as mechanism, not results.

\subsubsection{A hull test depended on the triangulation
(NR-23)}\label{a-hull-test-depended-on-the-triangulation-nr-23}

After one arm bought one CFD point, 6 of its next 20 designs were
rejected as extrapolations. A convex hull cannot shrink when a point is
added. A frozen variable written as a length and divided back gave a
unit coordinate of 1 + 2.2×10\textsuperscript{−16}, exactly on a facet,
and membership depended on which triangulation was built: 48 of 48 query
points inside the 56-point hull, 1 of 48 inside the 57-point hull. In
the adaptive-fidelity study this would have punished every arm that buys
CFD. It was found because an aborted run's log was read before being set
aside.

\subsubsection{A convergence pass nearly admitted two drifting CFD cases
(NR-27)}\label{a-convergence-pass-nearly-admitted-two-drifting-cfd-cases-nr-27}

The restart code stopped as soon as the criterion was met after two
blocks. Two cases met it in exactly one window that followed windows
which had not; one's window-mean C\_D alternated by 0.45\% between
blocks. The rule now requires two consecutive passing blocks. It was
written during the run after an interim look, which is disclosed; it is
stricter, and it removed two points the looser reading would have put
into the surface.

\subsubsection{Generated reports carried sentences from an older model
(NR-31,
NR-34)}\label{generated-reports-carried-sentences-from-an-older-model-nr-31-nr-34}

The optimiser-comparison report labelled 4038 drag-surface look-ups as
``CFD calls'' in a study that ran the solver zero times. The uncertainty
report printed a projected throughput as a measured one, understated its
sampling error about tenfold, and drew an attribution figure whose
y-axis labels belonged to another panel: the 0.609 bar labelled as entry
angle was the Sutton--Graves constant. Every table number matched its
JSON. The generators were corrected once no study held the source tree,
and a checker confirmed that no stored value moved: 2,551 and 13,744
prior values unchanged, and 272 of 275 table rows identical with 3
differing for declared reasons
(\texttt{REPORT\_\allowbreak{}REGENERATION\_\allowbreak{}2026-\allowbreak{}09-\allowbreak{}21.md}).

\subsection{The failure that was organisational (NR-18, NR-24)}

A two-hour optimiser-comparison study was launched as a detached job.
While it ran, a concurrent work-stream changed the evaluator's physics.
The runner had snapshotted its configuration, but worker processes
import the package when spawned, so an edit to the source can split one
candidate log between two physics models with nothing in the log to say
which row is which. The header said only ``dirty tree''. The run was
killed; no summary was written and no number from it was inspected; the
directory is kept with a note that it must not be analysed.

A configuration snapshot, a commit hash and a dirty flag are not
provenance for a long run in a repository several processes edit at
once. Every long study now hashes each source file of the physics
package at launch, re-checks after every logged batch, and on a mismatch
logs what was paid for, writes an abort marker and stops. A study
refuses to start on a screening made for a different design space.
Report generators no longer contain pre-written findings. The guard then
aborted two adaptive-fidelity development runs, correctly, because of
another work-stream's edits. This failure would have produced a
complete, plausible, well-formatted report.

\subsection{A gate that first read LIMITED (NR-25)}

§8.2 gives the account. The point for this section is that the
\texttt{LIMITED} assessment is kept on disk beside the \texttt{PASS},
and that the mean of the limit cycle was not the fixed point: the
converged C\_D differs from the oscillating window mean by about 0.03\%,
a third of the fine-to-medium difference, which moved the observed order
from 0.67 and 0.60 to 1.13 and 0.88.

\begin{figure}[htbp]
\centering
\includegraphics[width=\linewidth,height=0.42\textheight,keepaspectratio]{/Users/akj/Projects/aether/reports/final/../figures/infographics/07_gate_G4_history.pdf}
\caption{Gate G4 history: \texttt{LIMITED} to \texttt{PASS} by lowering the
maximum Courant number from 0.2 to 0.1 on the fine-mesh cases; the
restart rule was written after the first result had been seen and is
disclosed (NR-25, §8.2). Explanatory figure read from
\texttt{results/\allowbreak{}M2/\allowbreak{}M2-\allowbreak{}20260920T123901Z/\allowbreak{}gate\_\allowbreak{}assessment.json},
\texttt{gate\_\allowbreak{}assessment\_\allowbreak{}20260921\_\allowbreak{}0057\_\allowbreak{}before\_\allowbreak{}restarts.json},
\texttt{gci.csv},
\texttt{gci\_\allowbreak{}20260921\_\allowbreak{}0057\_\allowbreak{}before\_\allowbreak{}restarts.csv}
and \texttt{config\_\allowbreak{}snapshot.yaml}. File
\texttt{reports/\allowbreak{}figures/\allowbreak{}infographics/\allowbreak{}07\_\allowbreak{}gate\_\allowbreak{}G4\_\allowbreak{}history.png}.}
\end{figure}

\subsection{The experiment's traps, found before anything was built
(NR-11)}

Four traps surfaced on synthetic data, each of which would have been
invisible in a real experiment because the calibration would have fitted
well with wrong parameters. Contact resistance, fitted freely under a
prescribed flux, came back 186 to 531\% wrong and dragged conductivity
and specific heat 5 to 12\% low. A calibration step of about one
diffusion time left conductivity and specific heat each about 12\% low
while their ratio was recovered to 0.4\%. The maximum of a noisy record
is biased high, and averaging the near-peak plateau left +1.36 K of bias
because the window is selected by the noisy maximum; smoothing first cut
it to +0.08 K. And the ``seeded'' synthetic data used Python's salted
string hash, so it differed between processes while looking
deterministic within one.

\subsection{H2, and a study sized against an unreachable target (NR-35)}

§12 gives the result. The lesson recorded is procedural: before fixing a
hypervolume target, run the no-CFD arm and the reference search at the
planned budgets and check that the target can be reached.

\subsection{A validation exercise that changed nothing (NR-12)}

Three long-standing citations were opened for the first time, and not
one number changed (§6.3, §6.5). The same pass found that a deceleration
limit nobody had questioned corresponds to no document.

\subsection{Why these are kept}

Several numbers elsewhere depend on them: the Fidelity-0 fronts are
superseded because of §17.1, the TPI is absent from the objectives
because of §17.2, and the front of §10 has to be read through its
fences. In most of these cases a plausible, well-formatted, wrong answer
was available, and what caught it was not the test suite. It was a
figure with real axes, an interval outside its possible range, a budget
that went somewhere unexpected, a log that was read, or a guard that
fired when it was inconvenient.

\section{Conclusion}

Within a reduced-order model whose conduction solver, trajectory
integrator and heating constant were each checked against an external
reference, peak external heat flux and peak bondline temperature are
minimised by opposite ends of the tested range of entry angle. H0 is
supported in that domain.

On the final model, with drag from a CFD-derived surface and heating
corrected for the stagnation velocity gradient, the knee of the
two-objective front runs 14.9 K cooler at the bondline than the
peak-flux-only optimum for 21.5 kW/m\textsuperscript{2} more peak flux,
and the paired difference under the declared uncertainties is −14.31 K,
negative in every draw and every epistemic branch. H1 is supported,
narrowly: it is a statement about entry steepness at a geometry set by
an unsourced mass-fraction limit, the edge of where CFD was run, a box
bound and a frozen shoulder. The knee's heat shield is 346 kg of a 350
kg vehicle. The nominal optima violate a non-thermal constraint in 29.6
to 46.6\% of draws; a chance-constrained knee reduces that to 3.20\% for
3.1\% more peak flux.

H2 was not supported. Adaptive use of CFD reached the pre-declared
target in none of five seeds, the target was unreachable at the study's
budget, extra CFD changed the no-CFD arm's score by 0.0001, and the
AI-guided arm was not run. A language-model agent beat the best
conventional optimiser at 50 and 100 evaluations and showed no measured
difference at 200 under a rule declared in advance; its early lead
cannot be separated from prior knowledge.

\begin{figure}[htbp]
\centering
\includegraphics[width=\linewidth,height=0.42\textheight,keepaspectratio]{/Users/akj/Projects/aether/reports/final/../figures/infographics/04_hypothesis_scorecard.pdf}
\caption{Hypothesis scorecard: H0 supported in the tested domain; H1 supported
narrowly; H2 not supported; the optimiser comparison of §11 mixed.
Explanatory figure; the caveat text paraphrases §4 of this paper and
every numeral is read from a file:
\texttt{results/\allowbreak{}M1/\allowbreak{}M1-\allowbreak{}20260902T130547Z/\allowbreak{}candidates.csv},
\texttt{results/\allowbreak{}M4/\allowbreak{}M4-\allowbreak{}OPT-\allowbreak{}20260920T205252Z/\allowbreak{}summary.json},
\texttt{results/\allowbreak{}M7/\allowbreak{}M7-\allowbreak{}UQ-\allowbreak{}20260920T233048Z/\allowbreak{}paired\_\allowbreak{}difference.json},
\texttt{results/\allowbreak{}M6/\allowbreak{}M6-\allowbreak{}AF-\allowbreak{}20260920T211148Z/\allowbreak{}summary.json},
\texttt{results/\allowbreak{}M5/\allowbreak{}M5-\allowbreak{}ABL-\allowbreak{}20260920T211134Z/\allowbreak{}summary.json},
\texttt{reports/\allowbreak{}milestones/\allowbreak{}M6\_\allowbreak{}adaptive\_\allowbreak{}fidelity.md}.
File
\texttt{reports/\allowbreak{}figures/\allowbreak{}infographics/\allowbreak{}04\_\allowbreak{}hypothesis\_\allowbreak{}scorecard.png}.}
\end{figure}

Nothing here is established about a real vehicle or material. The coupon
experiment has not been run and no external review has taken place.

A design chosen on peak heat flux alone is, in this model, the design
with the hottest bondline on the front.

\section{Reproducibility}

\subsection{Environment and commands}

Python 3.12; numpy, scipy, matplotlib, pyyaml, pandas, pyarrow. OpenFOAM
is needed only for the CFD stages, and the version used was v2512. Full
instructions and measured runtimes are in \texttt{REPRODUCIBILITY.md}.

\begin{center}
\small\centering
\begin{tabular}{@{}>{\raggedright\arraybackslash}p{\dimexpr(\linewidth-6\tabcolsep)*333/1000\relax}>{\raggedright\arraybackslash}p{\dimexpr(\linewidth-6\tabcolsep)*333/1000\relax}>{\raggedright\arraybackslash}p{\dimexpr(\linewidth-6\tabcolsep)*333/1000\relax}@{}}
\toprule
\tabhead{Command} & \tabhead{Produces} & \tabhead{Measured wall time} \\
\midrule
\texttt{make test} & verification suite (497 passed, 1 skipped at the 2026-09-21
regeneration) & see note below \\
\texttt{make burn-\allowbreak{}vs-\allowbreak{}bake} & §7 & about 4 min \\
\texttt{make cfd-\allowbreak{}validate} & §8, gate G4 (OpenFOAM) & hours; fine cases 14,247 s and 8,911 s on a loaded machine \\
\texttt{make cfd-\allowbreak{}design-\allowbreak{}points},
\texttt{make aero-\allowbreak{}surface},
\texttt{make m3-\allowbreak{}coupled} & §9 (OpenFOAM for the first) & about 55 min + 27 min for the v2 extension; seconds; about 1 min \\
\texttt{make doe \&\& make optimize} & §10 & 3 min 20 s; 7 min 23 s on 6 workers \\
\texttt{make ablation-\allowbreak{}replay RUN\_\allowbreak{}ID=M5-\allowbreak{}ABL-\allowbreak{}20260920T211134Z} & §11, with no LLM access & a few minutes \\
\texttt{make adaptive} & §12 (OpenFOAM) & 3 h 44 min \\
\texttt{make robust}, \texttt{make uncertainty} & §13 & 2 h 18 min; 25.0 min on 4 workers \\
\bottomrule
\end{tabular}
\end{center}

Wall times were measured on one fanless laptop, mostly while other
studies shared it. \texttt{REPRODUCIBILITY.md}'s \texttt{make test} line
was corrected in the 2026-09-21 consistency sweep to match the 498-test
suite (497 passed, 1 skipped) above; it previously described an earlier
38-test count.

\subsection{Determinism and provenance}

The Fidelity-0 chain contains no random sampling; repeated runs on one
commit give bit-identical candidate files. Every stochastic study takes
and records explicit seeds. Every run writes an immutable configuration
snapshot with run ID, config hash, git commit, dirty flag and timestamp,
and an append-only candidate log that includes infeasible and invalid
designs. Milestone reports are generated from result files, and a hand
edit is overwritten by the next run. The three hand-written companions
(the agent audit and the two post-hoc addenda) are labelled as such and
override no number. Aborted and rejected runs ship with the repository.

The language-model study is reproducible as a run, not as a model. Every
prompt and raw response is persisted, and a strict replay with no model
access reproduced every method's per-seed hypervolume exactly with zero
prompt mismatches, including after the report-generator edits, under a
different source hash which the replay records. Issuing new calls will
give different numbers.

The repository is not yet public or archived at a clean commit. The
paper was drafted against \texttt{de2a834}.

\subsection{AI usage}

Most of the software was written by Claude through Claude Code, working
from the project specification: the physics models, geometry, CFD
pipeline, optimisation, surrogate, adaptive-fidelity and uncertainty
machinery, test suite, and figure and report generators. This paper was
also drafted by the assistant from the result files.
\texttt{AI\_\allowbreak{}USAGE.md} records each component in one of four
categories.

Rows marked \emph{review pending} mean that I have not yet read that
code line by line. As of the last measurement recorded in
\texttt{AI\_\allowbreak{}USAGE.md} the package was 23,869 lines of
Python with 6,587 lines of tests, nearly all added in one day, and
almost none of it read. A reading order of about 16 hours is in
\texttt{docs/\allowbreak{}ai\_\allowbreak{}usage\_\allowbreak{}proposed\_\allowbreak{}update.md}.
Twenty-two modelling decisions are recorded in
\texttt{AI\_\allowbreak{}USAGE.md} as AI-proposed with my approval
pending, and none has been approved. Where this paper says ``I did not
change the constant'' or ``I have left the criterion as written'', the
decision was proposed by the assistant and has stood; formal approval is
one of the things that remains for me to do.

Defects introduced by generated code are documented: a divergent
boundary-condition solve (caught by the solver failing), a truncated
thermal solve (caught by physical reasoning about a number that looked
low), an inverted dominance test (caught by a figure), an estimator with
impossible intervals (caught by the interval), a surrogate that returned
NaN (caught by where the budget went), and report sentences written for
an older model (caught by reading reports against their own data and
figures).

The assistant did not contribute the research question, the hypotheses,
the frozen scope, or the choice of the bondline as the failure mode of
interest. Those came from the specification.

\FloatBarrier

\section{References}

Tiers follow the project's scale: \tierok{} T1 opened and read in the
primary document; \tiermid{} T2 reputable secondary, or primary reached
only in part. How each was read is recorded in
\texttt{docs/\allowbreak{}validation/\allowbreak{}sourcing\_\allowbreak{}report.md},
\texttt{docs/\allowbreak{}theory/\allowbreak{}*} or the data file named.

\subsection*{Verified by this project}\label{verified-by-this-project}
\addcontentsline{toc}{subsection}{Verified by this project}

\begin{enumerate}
\def\labelenumi{\arabic{enumi}.}
\tightlist
\item
  \tierok{} T1. Sutton, K. and Graves, R. A. Jr., \emph{A General
  Stagnation-Point Convective-Heating Equation for Arbitrary Gas
  Mixtures}, NASA TR R-376, 1971. NTRS 19720003329.
\item
  \tiermid{} T2. NASA TFAWS 2012, \emph{Aerothermodynamics Course},
  Lecture 1. NASA-authored training material, not peer reviewed; opened
  and read.
\item
  \tiermid{} T2. Carroll, B. and Brandis, A., ``Stagnation Point
  Convective Heating Correlations for Entry into H\textsubscript{2}/He
  Atmospheres'', AIAA Aviation 2022. NTRS 20220018610.
\item
  \tiermid{} T2. ONERA/CNES, ``Noncatalytic and Finite Catalytic Heating
  Models for Atmospheric Re-entry Codes'', 1st International Orbital
  Debris Conference, 2019.
\item
  \tierok{} T1. Zoby, E. V. and Sullivan, E. M., \emph{Effects of Corner
  Radius on Stagnation-Point Velocity Gradients on Blunt Axisymmetric
  Bodies}, NASA TM X-1067, 1965. NTRS 19660017753.
\item
  \tierok{} T1. Ellison, J. C., \emph{Experimental Stagnation-Point
  Velocity Gradients and Heat-Transfer Coefficients for a Family of
  Blunt Bodies at Mach 8 and Angles of Attack}, NASA TN D-5121, 1969.
  NTRS 19690013192.
\item
  \tierok{} T1. Zoby, E. V., \emph{Empirical Stagnation-Point
  Heat-Transfer Relation in Several Gas Mixtures at High Enthalpy
  Levels}, NASA TN D-4799, 1968. NTRS 19680025996.
\item
  \tierok{} T1. Allen, H. J. and Eggers, A. J. Jr., NACA Report 1381,
  1958. NTRS 19930091020.
\item
  \tierok{} T1. Putnam, Z. R. and Braun, R. D., ``Extension and
  Enhancement of the Allen--Eggers Analytical Ballistic Entry Trajectory
  Solution'', \emph{JGCD} 38(3), 2015. DOI 10.2514/1.G000846.
\item
  \tierok{} T1. Desai, P. N. and Qualls, G. D., ``Stardust Entry
  Reconstruction'', AIAA 2008-1198. NTRS 20080008567.
\item
  \tierok{} T1. Pavlosky, J. E. and St.~Leger, L. G., \emph{Apollo
  Experience Report: Thermal Protection Subsystem}, NASA TN D-7564,
  1974. NTRS 19740007423.
\item
  \tierok{} T1. Smith, Hamermesh and Hogenson, NASA CR-159900, 1979.
  NTRS 19790012947.
\item
  \tierok{} T1. Vander Kam and Gage, ``Estimating Orion Heat Shield
  Failure Due to Ablator Cracking During the EFT-1 Mission''. DOI
  10.2514/1.A35003; NTRS 20160007549.
\item
  \tierok{} T1. NASA-STD-3001 Vol 2 Rev F, Table 6.5-1.
\item
  \tierok{} T1. NASA TM-104753, 1992. NTRS 19930020462.
\item
  \tiermid{} T2. Martin et al., ``ITRUSST Consensus on Standardised
  Reporting for Transcranial Ultrasound Stimulation'', arXiv:2402.10027,
  2024. Opened; reproduces the CEM43 formula. The Sapareto and Dewey
  original was not opened.
\item
  \tiermid{} T2. Ye, H. and De, S., ``Thermal injury of skin and
  subcutaneous tissues'', 2016. PMC5459687. Opened; the Henriques
  originals were not.
\item
  \tierok{} T1. Chapman, D. R., NACA Report 1051, 1951. NTRS
  19930090963.
\item
  \tierok{} T1. Love, E. S., NACA TN 3819, 1957. NTRS 19930084517.
\item
  \tierok{} T1. Miller, C. G. III, NASA TN D-4800, 1968. NTRS
  19680026196.
\item
  \tierok{} T1. Hoerner, S. F., \emph{Fluid-Dynamic Drag}, 1965, Ch.
  XVI.
\item
  \tiermid{} T2. Karlgaard, Schoenenberger and Van Norman, ``Base Drag
  Model Derived from Mars 2020 Backshell Pressure Measurements'',
  \emph{JSR} 60(6), 2023. DOI 10.2514/1.A35774. Paywalled; abstract
  level only.
\item
  \tierok{} T1. Capsule geometry set: NASA/TM-2005-213457 (Apollo, NTRS
  20060008942); Erb and Jacobs (Mercury); Wilmoth, Mitcheltree and Moss,
  AIAA 97-2510 (Stardust); NTRS 20060017066 (Viking); Spencer et al.,
  AAS 98-146 and NTRS 20040090463 (Pathfinder); NTRS 20120011936,
  20090024218 (MSL); NTRS 20160000307 (Hayabusa).
\item
  \tierok{} T1. \emph{U.S. Standard Atmosphere, 1976}, NOAA-S/T 76-1562
  / NASA-TM-X-74335, Table I pp.~68--69.
\item
  \tierok{} T1. NASA Glenn NPARC Alliance, ``Examining Spatial (Grid)
  Convergence'', CFD verification and validation tutorial,
  grc.nasa.gov/\allowbreak{}www/\allowbreak{}wind/\allowbreak{}valid/\allowbreak{}tutorial/\allowbreak{}spatconv.html.
  Source of the GCI formulae as used.
\item
  Opened and read; recorded in
  \texttt{data/\allowbreak{}reference/\allowbreak{}sphere\_\allowbreak{}supersonic.yaml},
  not in the sourcing report's graded list. Van Dyke, M. D. and Gordon,
  H. D., \emph{Supersonic Flow Past a Family of Blunt Axisymmetric
  Bodies}, NASA TR R-1, 1959, Table III. NTRS 19980223623.
\item
  Opened and read; recorded in the same data file. Bailey, A. B. and
  Hiatt, J., AEDC-TR-70-291, 1971 (DTIC AD0721208): Table II (sphere
  total drag) and Sec. 4.5 (Clark's forebody pressure-drag expression as
  quoted there; Clark's and Hoerner's originals were not fetched).
\item
  Tiered \tierok{} T1 in \texttt{configs/\allowbreak{}uncertainty.yaml}
  as reported by the M7 milestone report: NASA-TM-4715, Tables 3.1 and
  3.3; NASA/TM-20210022157, pp.~63--65 (density standard deviations
  below 86 km).
\end{enumerate}

\subsection*{Cited but not verified in primary by this
project}\label{cited-but-not-verified-in-primary-by-this-project}
\addcontentsline{toc}{subsection}{Cited but not verified in primary by
this project}

These are used by the code or named in generated reports. None has been
checked against the primary text. Each should be verified before
submission, or the dependent statement reworded.

\begin{itemize}
\tightlist
\item
  Celik, I. B. et al., ``Procedure for Estimation and Reporting of
  Uncertainty Due to Discretization in CFD Applications'', \emph{J.
  Fluids Eng.} 130(7), 078001, 2008. DOI 10.1115/1.2960953.
  Bibliographic record verified via Crossref; content not opened. The
  sign term in the apparent-order iteration and the
  oscillatory-convergence rule in \texttt{cfd/\allowbreak{}gci.py} are
  unverified against any opened text.
\item
  Roache, P. J., the GCI and its safety factors. Nothing opened;
  citation details are from memory. Not cited directly; {[}25{]} is
  cited instead.
\item
  Billig, F. S., ``Shock-Wave Shapes around Spherical- and
  Cylindrical-Nosed Bodies'', \emph{J. Spacecraft Rockets} 4(6),
  822--823, 1967. DOI 10.2514/3.28969. Record verified; formula
  confirmed from secondary sources, one of which prints the exponent as
  3.2 where two print 3.24. Used for domain sizing and one comparison
  that decides nothing in G4.
\item
  Carslaw, H. S. and Jaeger, J. C., \emph{Conduction of Heat in Solids},
  2nd ed., §2.9 (semi-infinite solid under constant surface flux). The
  solution is standard and the agreement is 0.002\%, but the edition,
  section and printed form have not been checked against a copy. Gate G3
  is measured against it.
\item
  Saltelli (2010) first-order and Jansen total-order Sobol' estimators.
  No provenance record. The implementation is pinned against the
  Ishigami function's analytical indices.
\item
  Tauber, M., NASA TP-2914, 1989; Tauber and Sutton, 1991. Provenance is
  contradictory inside this repository: one engineering-notebook entry
  lists TP-2914 among primaries obtained and read, while the sourcing
  report and the validation matrix record it as not opened (\tierno{}
  T3). No claim in this paper rests on it.
\item
  Anderson, J. D., \emph{Hypersonic and High-Temperature Gas Dynamics},
  2nd ed., AIAA, 2006, §4.3 ``Mach-Number Independence''. Chapter and
  section titles verified against the Library of Congress table of
  contents; the text was not opened.
\item
  NACA Report 1135, eq. 100 (Rayleigh pitot relation). Used as a closed
  form; no record that the report was opened.
\item
  Rodrigues, ``Closed-Form Reconstruction of Zoby--Sullivan
  Stagnation-Point Heat-Flux Scaling'', \emph{JTHT}, 2026. DOI
  10.2514/1.T7458. Paywalled, not opened (\tierno{} T3).
\item
  Vinh, Busemann and Culp (1980); Regan and Anandakrishnan (1993);
  Cooper and Holloway, ``The Shuttle Tile Story'' (1981). Not obtainable
  (\tierno{} T3); named as gaps only.
\item
  NSGA-II and ParEGO are named as algorithms (the former through the
  pymoo library). No primary source for either was opened.
\end{itemize}

\clearpage\appendix\renewcommand{\thefigure}{A.\arabic{figure}}\setcounter{figure}{0}

\section{CFD flow-field gallery}

Generated by
\texttt{scripts/\allowbreak{}render\_\allowbreak{}cfd\_\allowbreak{}gallery.py}
(\texttt{make cfd-\allowbreak{}gallery}) from the case directories under
\texttt{cfd/\allowbreak{}generated/\allowbreak{}} and the result tables
under \texttt{results/\allowbreak{}}. Every figure is written as PNG and
vector PDF. Colour scales are perceptually uniform (viridis / magma /
cividis) and stated on each colour bar. All fields are drawn in the
meridional plane (x, r), axis equal, one flat-shaded polygon per cell of
the wedge mesh (no interpolation). ``Time step'' is the saved iteration
of the local-time-stepping pseudo-time march. The index of all twelve
gallery figures, with cases and saved iterations, is
\texttt{reports/\allowbreak{}figures/\allowbreak{}cfd/\allowbreak{}INDEX.md};
four of them appear in the body (Figures 6, 7, 8 and 10). Within its
tested assumptions the CFD is used for forebody pressure drag, surface
pressure and perfect-gas shock shape only (§8.1).

\begin{figure}[htbp]
\centering
\includegraphics[width=\linewidth,height=0.28\textheight,keepaspectratio]{/Users/akj/Projects/aether/reports/final/../figures/cfd/M2_sphere_cp_meshes.pdf}
\caption{Run \texttt{M2-\allowbreak{}20260920T123901Z}, cases
\texttt{sphere\_\allowbreak{}M3\_\allowbreak{}coarse},
\texttt{sphere\_\allowbreak{}M3\_\allowbreak{}medium},
\texttt{sphere\_\allowbreak{}M3\_\allowbreak{}fine\_\allowbreak{}Co0p1},
\texttt{sphere\_\allowbreak{}M6\_\allowbreak{}coarse},
\texttt{sphere\_\allowbreak{}M6\_\allowbreak{}medium},
\texttt{sphere\_\allowbreak{}M6\_\allowbreak{}fine\_\allowbreak{}Co0p1},
sphere R = 0.5 m, three meshes (coarse ×1, medium ×2, fine ×4 = the
max-Courant-0.1 restart of record), surface pressure sampled at each
case's final iteration (15,000, 20,000, 110,000 at Mach 3; 10,000,
20,000, 110,000 at Mach 6). rhoCentralFoam, axisymmetric Euler,
calorically perfect gas γ = 1.4, p∞ = 1000 Pa, T∞ = 220 K. C\_p = (p −
p∞)/q∞. Dashed: modified Newtonian theory C\_p = C\_p,max
cos\textsuperscript{2}θ with C\_p,max from the exact Rayleigh-pitot
relation, an analytical estimate for comparison, not a validation
reference. Line style and colour distinguish the meshes. File
\texttt{reports/\allowbreak{}figures/\allowbreak{}cfd/\allowbreak{}M2\_\allowbreak{}sphere\_\allowbreak{}cp\_\allowbreak{}meshes.png}.}
\end{figure}

\begin{figure}[htbp]
\centering
\includegraphics[width=\linewidth,height=0.28\textheight,keepaspectratio]{/Users/akj/Projects/aether/reports/final/../figures/cfd/M2_demo_capsule_M6_fields.pdf}
\caption{Run \texttt{M2-\allowbreak{}20260920T123901Z}, case
\texttt{capsuledemo\_\allowbreak{}M6\_\allowbreak{}x2\_\allowbreak{}a1}:
a generic blunted 60° cone (nose R 0.6 m, D 1.2 m; not a flown vehicle,
not a validated result), M∞ = 6, medium mesh (refinement factor 2,
12,288 cells, sizing radius 1 m), saved iteration 30,000: Mach number,
p/p∞ and T/T∞. rhoCentralFoam, axisymmetric Euler, calorically perfect
gas γ = 1.4, p∞ = 1000 Pa, T∞ = 220 K. White: sonic line. Red: shock
stand-off Δ from \texttt{case\_\allowbreak{}result.json}; blue: distance
from the captured shock to the fixed-value inflow boundary on the axis
(x\_inflow = −0.313 m), the clearance NR-26 is about: the first attempt
in the sphere-sized domain,
\texttt{capsuledemo\_\allowbreak{}M6\_\allowbreak{}x2} (sizing radius
0.6 m), failed in the solver and wrote no field. C\_D,fore = 1.3967.
Colour: viridis / magma / cividis. File
\texttt{reports/\allowbreak{}figures/\allowbreak{}cfd/\allowbreak{}M2\_\allowbreak{}demo\_\allowbreak{}capsule\_\allowbreak{}M6\_\allowbreak{}fields.png}.}
\end{figure}

\begin{figure}[htbp]
\centering
\includegraphics[width=\linewidth,height=0.34\textheight,keepaspectratio]{/Users/akj/Projects/aether/reports/final/../figures/cfd/M3_design_point_cp.pdf}
\caption{Run \texttt{M3-\allowbreak{}DP-\allowbreak{}20260920T1610Z}, cases
\texttt{sc1p2\_\allowbreak{}coarse\_\allowbreak{}a0},
\texttt{dp017\_\allowbreak{}coarse\_\allowbreak{}a0},
\texttt{dp057\_\allowbreak{}coarse\_\allowbreak{}a0},
\texttt{dp001\_\allowbreak{}coarse\_\allowbreak{}Co0p05},
\texttt{dp007\_\allowbreak{}coarse\_\allowbreak{}a0},
\texttt{dp011\_\allowbreak{}coarse\_\allowbreak{}Co0p05},
\texttt{dp061\_\allowbreak{}coarse\_\allowbreak{}Co0p05}, coarse mesh
(3,072 cells), D = 1.2 m, surface pressure sampled at each case's final
iteration (15,000, 10,000, 20,000, 105,000, 15,000, 135,000, 135,000).
rhoCentralFoam, axisymmetric Euler, calorically perfect gas γ = 1.4, p∞
= 1000 Pa, T∞ = 220 K. C\_p = (p − p∞)/q∞ against arc length from the
nose over the diameter; the left panels show the same outlines in the
same colour and line style. Series tagged \texttt{[REJECTED]} are the
two NR-27 shapes: their C\_p is the last Courant restart's final save
and is not on the drag surface. At Mach 6 only the baseline shape
(scale-check case \texttt{sc1p2}) exists. File
\texttt{reports/\allowbreak{}figures/\allowbreak{}cfd/\allowbreak{}M3\_\allowbreak{}design\_\allowbreak{}point\_\allowbreak{}cp.png}.}
\end{figure}

\begin{figure}[htbp]
\centering
\includegraphics[width=\linewidth,height=0.34\textheight,keepaspectratio]{/Users/akj/Projects/aether/reports/final/../figures/cfd/M3_mach27_dp062_startup_failure.pdf}
\caption{Runs \texttt{M3-\allowbreak{}DP-\allowbreak{}20260920T1610Z} (top row:
\texttt{dp062\_\allowbreak{}coarse\_\allowbreak{}a2}, the design-point
attempt of record) and
\texttt{M3-\allowbreak{}mach27-\allowbreak{}startup-\allowbreak{}exp}
(bottom row: \texttt{dp062\_\allowbreak{}A\_\allowbreak{}fo6000}, the
NR-28 isolation run with a 6,000-iteration first-order start), point
\texttt{dp062}, M∞ = 27, coarse mesh (3,072 cells): T/T∞, Mach number
and kinetic-energy fraction of the first-order start-up solution.
rhoCentralFoam, axisymmetric Euler, calorically perfect gas γ = 1.4, p∞
= 1000 Pa, T∞ = 220 K. What survives is only the field written at the
end of the first-order (upwind) start-up, iteration 1,500 and 6,000
respectively. The van Leer hand-over then diverged with a floating-point
exception in sqrt (a negative temperature) at iteration 1,567 and 6,085
respectively, and a crashed run writes no field: the cell where the
temperature went negative is not recorded anywhere on disk and cannot be
shown. The saved first-order fields contain no cell below T∞ (minimum T
= 220 K), so they hold no precursor either. The third column shows the
mechanism NR-28 names: the kinetic energy is above 99\% of the total
energy in the freestream and stays dominant through the thin attached
shock layer along the 20° cone, so a small error in total energy is a
large relative error in c\_v T. Other attempts, all without a
post-hand-over field:
\texttt{dp062\_\allowbreak{}coarse\_\allowbreak{}a0} died at iteration
1,528; \texttt{dp062\_\allowbreak{}coarse\_\allowbreak{}a1} at 1,554;
\texttt{dp063\_\allowbreak{}coarse\_\allowbreak{}a2} at 1,601;
\texttt{dp062\_\allowbreak{}B\_\allowbreak{}fo1500\_\allowbreak{}co01}
at 1,569;
\texttt{dp062\_\allowbreak{}C\_\allowbreak{}fo6000\_\allowbreak{}co01}
at 6,047. Colour: cividis (T/T∞, energy fraction), viridis (Mach). File
\texttt{reports/\allowbreak{}figures/\allowbreak{}cfd/\allowbreak{}M3\_\allowbreak{}mach27\_\allowbreak{}dp062\_\allowbreak{}startup\_\allowbreak{}failure.png}.}
\end{figure}

\begin{figure}[htbp]
\centering
\includegraphics[width=\linewidth,height=0.66\textheight,keepaspectratio]{/Users/akj/Projects/aether/reports/final/../figures/cfd/M2_sphere_M3_fine_fourpanel.pdf}
\caption{Run \texttt{M2-\allowbreak{}20260920T123901Z}, case
\texttt{sphere\_\allowbreak{}M3\_\allowbreak{}fine\_\allowbreak{}Co0p1},
sphere R = 0.5 m, M∞ = 3, fine mesh (refinement factor 4, 49,152 cells),
max Courant 0.1, the solution of record after NR-25, saved iteration
110,000: Mach number, p/p∞, T/T∞ and ρ/ρ∞. rhoCentralFoam, axisymmetric
Euler, calorically perfect gas γ = 1.4, p∞ = 1000 Pa, T∞ = 220 K. One
flat polygon per cell. White line: sonic line M = 1 from the cell-centre
triangulation; it bounds the subsonic nose region and also runs along
the captured shock, where M passes through 1 inside the one to two cells
of numerical shock thickness. Dashed: Billig's correlation for the shock
shape (independent empirical curve, not a fit). Δ: shock stand-off from
\texttt{case\_\allowbreak{}result.json} (50\% density-rise point on the
stagnation line). Colour bars: Mach viridis, p/p∞ magma, T/T∞ cividis,
ρ/ρ∞ viridis. File
\texttt{reports/\allowbreak{}figures/\allowbreak{}cfd/\allowbreak{}M2\_\allowbreak{}sphere\_\allowbreak{}M3\_\allowbreak{}fine\_\allowbreak{}fourpanel.png}.}
\end{figure}

\begin{figure}[htbp]
\centering
\includegraphics[width=\linewidth,height=0.66\textheight,keepaspectratio]{/Users/akj/Projects/aether/reports/final/../figures/cfd/M2_sphere_M3_limit_cycle_field.pdf}
\caption{Run \texttt{M2-\allowbreak{}20260920T123901Z}, sphere R = 0.5 m, M∞ = 3,
fine mesh (49,152 cells), cases
\texttt{sphere\_\allowbreak{}M3\_\allowbreak{}fine},
\texttt{sphere\_\allowbreak{}M3\_\allowbreak{}fine\_\allowbreak{}cyclediag},
\texttt{sphere\_\allowbreak{}M3\_\allowbreak{}fine\_\allowbreak{}Co0p1}.
rhoCentralFoam, axisymmetric Euler, calorically perfect gas γ = 1.4, p∞
= 1000 Pa, T∞ = 220 K. (a) Per-cell standard deviation of p over the 100
snapshots of the cycle-diagnosis continuation of the max-Courant-0.2
case (written every 2 iterations), divided by the per-cell mean: the
fluctuation is smallest around the stagnation point and grows along the
body in ray-like bands from the captured shock toward the supersonic
outflow (NR-25). (b) The same case's difference between its two last
saved fields, a two-phase sample of the oscillation. (c) The
max-Courant-0.1 restart of record, difference between its two last saved
fields, on the same colour scale; the intervals differ (20,000 against
5,000 iterations), so read the level, not a ratio. The two-time
difference does not separate the two Courant numbers by its maximum
(5.0×10\textsuperscript{−2} against 2.0×10\textsuperscript{−2}, both in
the cells the captured shock straddles near the outflow corner); it does
by its level inside the shock layer: median 1.6×10\textsuperscript{−3}
against 2.4×10\textsuperscript{−4}, and 58\% against 19\% of the layer's
cells above 10\textsuperscript{−3}. ``Shock layer'' = cells with ρ/ρ∞
\textgreater{} 1.5 at the later saved iteration. (d) p/p∞ of record.
Colour: cividis, log\textsubscript{10} clipped to {[}−4.5, −1{]}; magma
for p/p∞. File
\texttt{reports/\allowbreak{}figures/\allowbreak{}cfd/\allowbreak{}M2\_\allowbreak{}sphere\_\allowbreak{}M3\_\allowbreak{}limit\_\allowbreak{}cycle\_\allowbreak{}field.png}.}
\end{figure}

\begin{figure}[htbp]
\centering
\includegraphics[width=\linewidth,height=0.66\textheight,keepaspectratio]{/Users/akj/Projects/aether/reports/final/../figures/cfd/M2_sphere_M6_limit_cycle_field.pdf}
\caption{As Figure A.6, for M∞ = 6: run
\texttt{M2-\allowbreak{}20260920T123901Z}, cases
\texttt{sphere\_\allowbreak{}M6\_\allowbreak{}fine},
\texttt{sphere\_\allowbreak{}M6\_\allowbreak{}fine\_\allowbreak{}cyclediag},
\texttt{sphere\_\allowbreak{}M6\_\allowbreak{}fine\_\allowbreak{}Co0p1},
fine mesh (49,152 cells). Same gas model and freestream. The two-time
difference does not separate the two Courant numbers by its maximum
(1.1×10\textsuperscript{−1} against 1.2×10\textsuperscript{−1}, both in
the cells the captured shock straddles near the outflow corner); it does
by its level inside the shock layer: median 3.6×10\textsuperscript{−3}
against 9.0×10\textsuperscript{−4}, and 76\% against 47\% of the layer's
cells above 10\textsuperscript{−3}. ``Shock layer'' = cells with ρ/ρ∞
\textgreater{} 1.5 at the later saved iteration. (d) p/p∞ of record.
Colour: cividis, log\textsubscript{10} clipped to {[}−4.5, −1{]}; magma
for p/p∞. File
\texttt{reports/\allowbreak{}figures/\allowbreak{}cfd/\allowbreak{}M2\_\allowbreak{}sphere\_\allowbreak{}M6\_\allowbreak{}limit\_\allowbreak{}cycle\_\allowbreak{}field.png}.}
\end{figure}

\begin{figure}[htbp]
\centering
\includegraphics[width=\linewidth,height=0.66\textheight,keepaspectratio]{/Users/akj/Projects/aether/reports/final/../figures/cfd/M6_adaptive_promotions_mach.pdf}
\caption{Run \texttt{M6-\allowbreak{}AF-\allowbreak{}20260920T211148Z}, adaptive
arm, coarse mesh (3,072 cells), D = 1.2 m, forebody domain, M∞ = 20,
each panel at its case's final saved iteration (title): Mach number.
rhoCentralFoam, axisymmetric Euler, calorically perfect gas γ = 1.4, p∞
= 1000 Pa, T∞ = 220 K. A promotion is a candidate the adaptive policy
sent to CFD instead of trusting the Fidelity-1 surface
(\texttt{promotion\_\allowbreak{}log.json} reason for every adaptive
promotion: ``promising, and the cheap model is untrusted here''; none
was outside the arm's hull). Every usable case is appended to that
arm-and-seed's own surface, which is re-fitted (\texttt{counters.json}
surface history):
\texttt{af2a27042848\_\allowbreak{}coarse\_\allowbreak{}a0}, seed 37,
usable, surface training set 69 {\symfont →} 70 {\symfont →} 71
{\symfont →} 72 {\symfont →} 74 points over the run;
\texttt{af5e13e211ab\_\allowbreak{}coarse\_\allowbreak{}a0}, seed 37,
usable, 69 {\symfont →} 70 {\symfont →} 71 {\symfont →} 72 {\symfont →}
74; \texttt{af01fbe28a96\_\allowbreak{}coarse\_\allowbreak{}a0}, seed
67, usable, 69 {\symfont →} 70 {\symfont →} 71 {\symfont →} 72;
\texttt{af1376216f0b\_\allowbreak{}coarse\_\allowbreak{}Co0p05}, seed
37, rejected, 69 {\symfont →} 70 {\symfont →} 71 {\symfont →} 72
{\symfont →} 74. The rejected case never met the force criterion in two
consecutive blocks at Courant 0.2, 0.1 and 0.05 (the NR-27 rule); its
field is the last Courant restart's final save, shown but not used.
White: sonic line. Colour: viridis. File
\texttt{reports/\allowbreak{}figures/\allowbreak{}cfd/\allowbreak{}M6\_\allowbreak{}adaptive\_\allowbreak{}promotions\_\allowbreak{}mach.png}.}
\end{figure}

\end{document}
